%--------------------------------------------------------------------------------------
% ÁLTALÁNOS ALAPSABLON A (RÉSZ)BEADANDÓKHOZ
%--------------------------------------------------------------------------------------
\documentclass[12pt, a4paper]{book}
\usepackage[margin=1in]{geometry}
\usepackage[utf8]{inputenc}
\usepackage[magyar]{babel}
\usepackage[hidelinks]{hyperref} 
\usepackage{graphicx}
\usepackage{comment}
\usepackage{amsthm}
\usepackage{MnSymbol}
\usepackage{wasysym}
\usepackage{fancyhdr}
\usepackage{longtable}
\usepackage{changepage}
\usepackage{xcolor}
\usepackage{algorithmicx}
\usepackage{listings}

\pagestyle{fancy}
\fancyhf{}
\fancyhead[LE,RO]{\thepage}
\fancyhead[RE,LO]{\leftmark}
%kitöltendő
\fancyfoot[CE,CO]{Csapatsorszám - Viseguard Studios}

\title{Szoftver projekt labor beadandó} 
%kitöltendő
\author{Karsa Zoltán, Killin Attila, Jakab Dániel, Kacsó Péter, Imets Tamás} 

\setcounter{secnumdepth}{4}

\begin{document}
\begin{titlepage} 
	\begin{center}
		\includegraphics[width=5cm]{docs/img/BMElogo}\\ 
		\vspace{1cm}
		\LARGE{\bfseries{Szoftver projekt laboratórium}}\\
		\vspace{0.5cm}
		%kitöltendő
		\Large\textsc{X. Összesített dokumentum}\\
		\vspace{1cm}
		\small{Csapat}\\
		%kitöltendő
		\Large{\bfseries{Csapatsorszám - Viseguard Studios}}\\
		\vspace{1cm}
		\small{Konzulens}\\
		%kitöltendő
		\Large{Dr. Goldschmidt Balázs}
	\end{center}
	\begin{flushleft}
		\vspace*{6cm}
		Csapattagok\\
		\vspace{-0.3cm}
		\rule{14cm}{0.5pt}\\
		\vspace{0.2cm}
		\begin{tabular}{l l l} %kitöltendő
			Sike Ádám & E8Z277 & molteneq.official@gmail.com \\
			Dömötör Péter & G2Y5TI & javadaniel@iit.bme.hu \\
			Projlab Péter & ?????? & projlabpeter@iit.bme.hu \\
			LaTex Zoltán  & ?????? & latexzoltan@iit.bme.hu \\
			Git Attila & ?????? & gitattila@iit.bme.hu \\
		\end{tabular}
	\end{flushleft}
	\begin{flushright}
		\vspace*{2cm}
		\today   %vagy saját dátum megadása, manuálisan
	\end{flushright}
\end{titlepage}
%---------------------------------------------------------------
%\\\\\\\\\\\\\\\\\\\\\\\\NE MÓDOSÍTSD///////////////////////////

\newenvironment{naplo}{\begin{longtable}{|l|l|p{2.5cm}|p{7.5cm}|}
        \hline
        \hline \multicolumn{1}{|c|}{\textbf{Kezdet}} & \multicolumn{1}{c|}{\textbf{Időtartam}} & \multicolumn{1}{c|}{\textbf{Résztvevők}} & \multicolumn{1}{c|}{\textbf{Leírás}}  \\ \hline 
        \endfirsthead
        
        \multicolumn{3}{c}%
        {{\bfseries}} \\
        \hline \multicolumn{1}{|c|}{\textbf{Kezdet}} & \multicolumn{1}{c|}{\textbf{Időtartam}} & \multicolumn{1}{c|}{\textbf{Résztvevők}} & \multicolumn{1}{c|}{\textbf{Leírás}}  \\ \hline 
        \endhead
        
}{\end{longtable}}

\newcommand{\naplotag}[4]{#1 & #2 & #3 & #4 \\
    \hline}

\newtheorem{defi}{Definíció}
\newtheorem{rov}{Rövidítés}
\newcommand{\urlref}[2]{
	\parbox{\linewidth}{\href{#1}{\RIGHTarrow \quad #2 - \texttt{#1}}}}

\newcommand{\diagram}[3]{
	\begin{figure}[h]
		\centering
		\includegraphics[width=#3]{#1}
		\caption{#2}
	\end{figure}}

\renewcommand{\comment}[1]{\begin{flushright}
		\color{red} 
		#1
\end{flushright} }

\newenvironment{class-template}[5]
{ 
    \hspace*{0.4em} $\blacksquare $ \large Felelősség \\  \normalsize #1 \\[12px]
    \hspace*{0.4em} $\blacksquare $ \large Ősosztályok \\  \normalsize #2 \\[12px]
    \hspace*{0.4em} $\blacksquare $ \large Interfészek \\  \normalsize #3 \\[12px]
    \hspace*{0.4em} $\blacksquare $ \large Attribútumok \\  \normalsize #4 \\[12px]
    \hspace*{0.4em} $\blacksquare $ \large Metódusok \\  \normalsize #5 \\[12px]
}{}

\newenvironment{class-template-responsibility}{\hspace*{0.4em} $\blacksquare $ \large Felelősség \\[5px]  \normalsize}{\\[12px]}
\newenvironment{class-template-baseclass}{\hspace*{0.4em} $\blacksquare $ \large Ősosztály \\[5px]  \normalsize}{\\[12px]}
\newenvironment{class-template-interface}{\hspace*{0.4em} $\blacksquare $ \large Interfészek \\[5px]  \normalsize}{\\[12px]}
\newenvironment{class-template-method}{\hspace*{0.4em} $\blacksquare $ \large Metódusok \\[5px]  \normalsize}{\\[12px]}
\newenvironment{class-template-attribute}{\hspace*{0.4em} $\blacksquare $ \large Attribútumok \\[5px]  \normalsize}{\\[12px]}
\newenvironment{class-template-statechart}{\hspace*{0.4em} $\blacksquare$ \large Állapotgép \\[5px] \normalsize}{\\[12px]}

\newenvironment{enum-template}[2]
{ 
    \hspace*{0.4em} $\blacksquare $ \large Felelősség \\[5px]  \normalsize #1 \\[12px]
    \hspace*{0.4em} $\blacksquare $ \large Literálok \\[5px]  \normalsize #2 \\[12px]
}{}	

{\vspace*{-1cm}}
\newcommand{\classitem}[2]{
    \hspace*{25px}\parbox{0.9\linewidth}{$\Diamond $ \textit{#1} - #2}
}

\newcommand{\baseclass}{ $\rightarrow$ }

\newenvironment{funkovetelmeny}[6]
{
    \noindent\begin{tabular}{| p{2cm} | p{2cm} | p{2.2cm} | p{5.3cm} | p{2.5cm} |}
        \hline
        \hline
        \textbf{Azonosító} & \textbf{Prioritás} & \textbf{Forrás} & \textbf{Use-case} & \textbf{Ellenőrzés} \\
        \hline
        #1 & #2 & #3 & #4 & #5 \\
        \hline
    \end{tabular}
    \\[1ex]
    \textbf{Leírás:} #6 \\[0.2ex]
    \itshape}
{\\[5px]}

\newenvironment{kovetelmeny}[5]
{ 
    \noindent\begin{tabular}{| p{2cm} | p{3.5cm} | p{3.5cm} | p{3.5cm} |}
        \hline
        \hline
        \textbf{Azonosító} & \textbf{Prioritás} & \textbf{Forrás} & \textbf{Ellenőrzés} \\
        \hline
        #1 & #2 & #3 & #4 \\
        \hline
    \end{tabular}
    \\[1ex]
    \textbf{Leírás:} #5 \\[0.2ex]
    \itshape}
{\\[5px]}

\newenvironment{use-case}[3]
{ 
    
    \vspace{10px}
    \noindent\begin{tabular}{| p{3cm} | p{12cm} |}
        \hline
        \hline
        \vspace{0.2ex} \textsc{\textbf{Use-case neve:}} & \vspace{0.2ex} \textsc{#1} \\[0.2ex]
        \hline
        \textbf{Rövid leírás:} & #2 \\
        \hline
        \textbf{Aktorok:} & #3 \\
        \hline
        \textbf{Forgatókönyv:} &  
    }
    {	\\
        \hline
\end{tabular} }

\newenvironment{fajllista}{
    \begin{longtable}{|p{4.5cm}|l|l| p{4.5cm} |}
        \hline
        \hline \multicolumn{1}{|c|}{\textbf{Fájl neve}} & \multicolumn{1}{c|}{\textbf{Méret}} & \multicolumn{1}{c|}{\textbf{Keletkezés ideje}} & \multicolumn{1}{c|}{\textbf{Tartalom}}  \\ \hline 
        \endfirsthead
        
        \multicolumn{3}{c}%
        {{\bfseries}} \\
        \hline \multicolumn{1}{|c|}{\textbf{Fájl neve}} & \multicolumn{1}{c|}{\textbf{Méret}} & \multicolumn{1}{c|}{\textbf{Keletkezés ideje}} & \multicolumn{1}{c|}{\textbf{Tartalom}}  \\ \hline 
        \endhead
        
    }{\end{longtable}}

\newcommand{\fajl}[4] {
    #1 & #2 & #3 & #4 \\ \hline
}

\newenvironment{ertekeles}{
    \begin{figure}[h]
        \centering
        \noindent\begin{tabular}{| p{5cm} | p{4cm} | p{4cm} |}
            \hline
            \hline
            \textbf{Tag neve} & \textbf{Tag neptun} & \textbf{Munka százalékban} \\	\hline
}{\end{tabular}\end{figure}}

\newcommand{\ertekelestag}[3]{ 
    #1 & #2 & #3 \\ \hline
}

\newcommand{\ertekelestagk}[2]{ 
    #1 & #2 \\ \hline
}

\newenvironment{ertekelesOra}{
    \begin{figure}[h]
        \centering
        \noindent\begin{tabular}{| p{5cm} | p{4cm} | p{4cm} |}
            \hline
            \hline
            \textbf{Tag neve} & \textbf{Tag neptun} & \textbf{Munka órában} \\	\hline
            
        }{\end{tabular}	\end{figure}}

\newenvironment{ertekelesKod}{
    \begin{figure}[h]
        \centering
        \noindent\begin{tabular}{| p{5cm} | p{4cm} |}
            \hline
            \hline
            \textbf{Fázis} & \textbf{Kódsorok száma} \\	\hline
            
        }{\end{tabular}	\end{figure}}

\newcommand{\test}[3]{
    \vspace{10px}
    \noindent\begin{tabular}{| p{3.5cm} | p{12cm} |}
        \hline
        \hline
        \vspace{0.2ex} \textsc{\textbf{Teszt-eset neve}} & \vspace{0.2ex} \textsc{#1} \\[0.2ex]
        \hline
        \textbf{Rövid leírás} & #2 \\
        \hline
        \textbf{Teszt célja} & #3 \\
        \hline
    \end{tabular}
}

\newenvironment{test-case}[4]
{ 
    \hspace*{0.4em} $\blacksquare $ \large Leírás   \\  \normalsize #1   \\[12px]
    \hspace*{0.4em} $\blacksquare $ \large Ellenőrzött funkcionalitás, várható hibahelyek  \\  \normalsize #2  \\[12px]
    \hspace*{0.4em} $\blacksquare $ \large Bemenet  \\  \normalsize #3  \\[12px]
    \hspace*{0.4em} $\blacksquare $ \large Elvárt kimenet \\  \normalsize #4 \\[12px]
}{}

\newenvironment{test-case-description}{\hspace*{0.4em} $\blacksquare $ \large Leírás \\[5px]  \normalsize}{\\[12px]}
\newenvironment{test-case-function}{\hspace*{0.4em} $\blacksquare $ \large Ellenőrzött funkcionalitás, várható hibahelyek \\[5px]  \normalsize}{\\[12px]}
\newenvironment{test-case-input}{\hspace*{0.4em}\parbox{\linewidth}{$\blacksquare $ \large Bemenet\vspace*{5px}}\normalsize}{\leavevmode\newline}
\newenvironment{test-case-output}{\hspace*{0.4em}\parbox{\linewidth}{$\blacksquare $ \large Elvárt kimenet\vspace{5px}}\normalsize}{\leavevmode\newline}

\newcommand{\testOK}[2]{
    \vspace{10px}
    \noindent\begin{tabular}{| p{4.2cm} | p{12cm} |}
        \hline
        \hline
        \textbf{Tesztelő neve} & \textsc{#1} \\[0.2ex]
        \hline
        \textbf{Teszt időpontja} & #2 \\
        \hline
    \end{tabular}
}

\newcommand{\testFAIL}[5]{
    \vspace{10px}
    \noindent\begin{tabular}{| p{4.2cm} | p{12cm} |}
        \hline
        \hline
        \textbf{Tesztelő neve} & \textsc{#1} \\[0.2ex]
        \hline
        \textbf{Teszt időpontja} & #2 \\
        \hline
        \textbf{Teszt eredménye} & #3 \\
        \hline
        \textbf{Lehetséges hibaokok} & #4 \\
        \hline
        \textbf{Változtatások} & #5 \\
        \hline
    \end{tabular}
}

\newenvironment{terv}{	
    \begin{longtable}{| l | p{8cm} | l |}
        \hline
        \hline
        \multicolumn{1}{|c|}{\textbf{Határidő}} & \multicolumn{1}{c|}{\textbf{Feladat}} & \multicolumn{1}{c|}{\textbf{Felelős}}  \\ \hline 
        \endfirsthead
        
        \multicolumn{3}{c}%
        {{\bfseries}} \\
        \hline \multicolumn{1}{|c|}{\textbf{Határidő}} & \multicolumn{1}{c|}{\textbf{Feladat}} & \multicolumn{1}{c|}{\textbf{Felelős}}  \\ \hline 
        \endhead
        
}{\end{longtable}}

\newcommand{\tervitem}[3]{ #1 & #2 & #3 \\ \hline}

\newenvironment{szotar}{
    \begin{longtable}{ p{4cm}  p{11cm} }
        
}{\end{longtable}}

\newcommand{\szotaritem}[2]{ \textbf{#1} & #2 \\[10px]}
\sloppy

%-----------------------------------------------------------------------------------------

%Végső beadáshoz tartalomjegyzék:
%\tableofcontents

%A fejezetek számozásának átállítására (mindig a megadott számtól folytatódik a számozás):
%\setcounter{chapter}{-1}


\setcounter{chapter}{1}
\chapter{Követelmény, projekt, funkcionalitás}

\section{Bevezetés}

\subsection{Cél}
%\comment{A dokumentum célja.}
\markdownInput{docs/bead02/feladat.md}

\subsection{Szakterület}
\comment{A kialakítandó szoftver milyen területen használható, milyen célra.}

\subsection{Definíciók, rövidítések}
\comment{A dokumentumban használt definíciók, rövidítések magyarázata.}
\begin{defi} 
	Példa defi: egy demó a defi használatára
\end{defi}
\begin{rov}
    PR: példa rövidítés
\end{rov}

\subsection{Hivatkozások}
\urlref{http://iit.bme.hu/}{BME IIT - Programozás alapjai 3. segédanyagok, Szoftvertechnikák segédanyagok, Szoftver projekt laboratórium feladatok}

\subsection{Összefoglalás}
\comment{A dokumentum további részeinek rövid ismertetése}

\section{Áttekintés}
\subsection{Általános áttekintés}
\comment{A kialakítandó szoftver legmagasabb szintű architekturális képe. A fontosabb alrendszerek felsorolása, a közöttük kialakítandó interfészek lényege, a felhasználói kapcsolatok alapja. Esetleges hálózati és adattárolási elvárások.}

\subsection{Funkciók}
\comment{A feladat kb. 4000 karakteres (kb 1,5 oldal) részletezettségű magyar nyelvű leírása. Nem szerepelhetnek informatikai kifejezések.}

\subsection{Felhasználók}
\comment{A felhasználók jellemzői, tulajdonságai}

\subsection{Korlátozások}
\comment{Az elkészítendő szoftverre vonatkozó – általában nem funkcionális - előírások, korlátozások.}

\subsection{Feltételezések, kapcsolatok}
\comment{A Hivatkozásokban felsoroltanyagok, web-oldalak kapcsolódása a feladathoz, melyik milyen szempontból érdekes, milyen inputot ad.}


\section{Követelmények}
\subsection{Funkcionális követelmények}
\comment{Az alábbi táblázat kitöltésével készítendő. Dolgozzon ki követelmény azonosító rendszert! Az ellenőrzés módja szokásosan bemutatás és/vagy kiértékelés. Prioritás lehet alapvető, fontos, opcionális. Az alapvető követelmények nem teljesítése végzetes. Forrás alatt a követelményt előíró anyagot, szervezetet kell érteni. Esetünkben forrás lehet maga a csapat is, mikor ő talál ki követelményt. Use-case-ek alatt az adott követelményt megvalósító használati esete(ke)t kell megadni.}
\begin{funkovetelmeny}
	%Azonosító
	{Azonosító}
	%Prioritás
	{Pioritás}
	%Forrás
	{Forrás}
	%Use case
	{Használati eset}
	%Ellenőrzés
	{llenőrzés}
	%Leírás
	{Leírás}
    Megjegyzés (opcionális)
\end{funkovetelmeny}


\subsection{Erőforrásokkal kapcsolatos követelmények}
\comment{A szoftver fejlesztésével és használatával kapcsolatos számítógépes, hardveres, alapszoftveres és egyéb architekturális és logisztikai követelmények}

\begin{kovetelmeny}
	{Azonosító} %Azonosító
	{Pioritás} %Prioritás
	{Forrás} %Forrás
	{Ellenőrzés} %Ellenőrzés
	{Leírás} %Leírás
    Megjegyzés (opcionális)
\end{kovetelmeny}

\subsection{Átadással kapcsolatos követelmények}
\comment{A szoftver átadásával, telepítésével, üzembe helyezésével kapcsolatos követelmények}
\begin{kovetelmeny}
    {Azonosító} %Azonosító
    {Pioritás} %Prioritás
    {Forrás} %Forrás
    {Ellenőrzés} %Ellenőrzés
    {Leírás} %Leírás
    Megjegyzés (opcionális)
\end{kovetelmeny}

\subsection{Egyéb nem funkcionális követelmények}
\comment{A biztonsággal, hordozhatósággal, megbízhatósággal, tesztelhetőséggel, a felhasználóval kapcsolatos követelmények}
\begin{kovetelmeny}
    {Azonosító} %Azonosító
    {Pioritás} %Prioritás
    {Forrás} %Forrás
    {Ellenőrzés} %Ellenőrzés
    {Leírás} %Leírás
    Megjegyzés (opcionális)
\end{kovetelmeny}

\section{Lényeges use-case-ek}
\comment{Funkcionális követelmények részben felsorolt követelmények közül az alapvető és fontos követelményekhez tartozó használati esetek megadása az alábbi táblázatos formában.}
\subsection{Use-case leírások}
\comment{Minden use-case-hez külön}
\begin{use-case}
	%név
	{Use-case Neve}
	%rövid leírás
	{Az eset rövid leírása}
	%aktorok
	{Aktorok}
	%forgatókönyv
	Forgatókönyv \newline 
        \textbf{A.1} Alternatíva
\end{use-case}

\subsection{Use-case diagram}
\diagram{docs/img/BMElogo}{Demó}{4cm}

\section{Szótár}
\comment{A szótár a követelmények alapján készítendő fejezet. Egy szótári bejegyzés definiálásához csak más szótári bejegyzések és köznapi – a feladattól független – fogalmak használhatók fel. A szótár mérete kb. 1-2 oldal legyen. A bejegyzések legyenek ABC sorrendben!}

%szotar.sh script, lsd. a dokumentációt!!!
\begin{szotar}
    \szotaritem{Kulcs 1 }{Érték 1}
    \szotaritem{Hóásás }{Egy réteg hó eltakarítása, azaz törlése az adott jégtábláról. Egy egység munkába kerül.}
\end{szotar}

\section{Projekt terv}
\comment{Tartalmaznia kell a projekt végrehajtásának lépéseit, a lépések, eredmények határidejét, az egyes feladatok elvégzéséért felelős személyek nevét és beosztását, a szükséges erőforrásokat, stb. Meg kell adni a csoportmunkát támogató eszközöket, a választott technikákat! Definiálni kell, hogy hogyan történik a dokumentumok és a forráskód megosztása!}


\subsection{Projektütemterv}
\begin{terv}
		\tervitem{febr. 4.}  {Követelmény, projekt, funkcionalitás}{ ?????? }
        \tervitem{márc. 2. } { Analízis modell kidolgozása 1. - beadás }{ ?????? } 
        \tervitem{márc. 9. } { Analízis modell kidolgozása 2. - beadás }{ ?????? }
        \tervitem{márc. 16. }{ Skeleton tervezése - beadás }{ ?????? }
        \tervitem{márc. 23. }{ Skeleton - beadás és a forráskód herculesre való feltöltése }{ ?????? }
        \tervitem{márc. 30. }{ Prototípus koncepciója - beadás }{ ?????? }
        \tervitem{ápr. 6. }{ Részletes tervek - beadás }{ ?????? }
        \tervitem{ápr. 27. }{ Prototípus - beadás és a forráskód, a tesztbemenetek és az elvárt kimenetek herculesre való feltöltése }{ ?????? }
        \tervitem{ máj. 4. }{ Grafikus felület specifikációja - beadás }{ ?????? }
        \tervitem{ máj. 18. }{ Grafikus változat és Összefoglalás - beadás és a forráskód herculesre való feltöltése }{ ?????? } 
\end{terv}

\subsection{Erőforrások, eszközök}
A fejlesztés során felhasznált segédeszközök:
\begin{itemize}
	\item Dokumentáció: dokumentáló eszközök felsorlása
	\item Kommunikáció: kommunikációs platformok...
	\item Modellező eszköz: modellezési eszközök
	\item Fejlesztő környezetek: ...
	\item Forráskód megosztás, verziókezelés: ...
	\item Egyéb ....
\end{itemize}

\comment{Még szabadon felvehető releváns idetartozó dolgok...}


\clearpage
\section{Napló}

\comment{ 
    A napló tartalmazza az előző beadás óta eltelt időszak történéseit időrendben. A naplóból egyértelműen ki kell derülnie, hogy az egyes anyagrészeket ki és mennyi idő alatt készítette.
    
    A napló bejegyzésekből áll. Minden bejegyzésnek tartalmaznia kell: 
    \begin{itemize}
        \item a történés kezdetének időpontját, nap-óra pontossággal
        \item történés időtartamát, óra felbontással
        \item a szereplő(k) nevét (Kérjük a szereplők VEZETÉKNEVÉT használni)
        \item a tevékenység leírását.
    \end{itemize}

    Amennyiben a tevékenységben több szereplő vesz részt, akkor aza tevékenység csak értekezlet lehet, amelynek az eredményei DÖNTÉSEK. A döntéseket precízen meg kell szövegezni (Pl.: Az X objektum Y és Z metódusainak kódját W készíti el Q határidőre).Ha a bejegyzés egyetlen személyhez kötődik, akkor meg kell adni, hogy a tevékenység milyen dologra irányul. A dolog a feladat kapcsán elkészítendő termék, amelynek a (esetleg korábban) beadott anyagban megtalálhatónak kell lenni.A naplóbejegyzés felbontásának egysége szöveges, rajzos anyag esetében az ábra, diagram, vagy kb. fél-egy oldalnyi szöveg. Kódban az egység a metódus. (Pl.: A 3. ábrán látható szekvencia-diagram kidolgozása, vagy az X objektum Y és Z metódusainak kódolása és belövése.
}

\begin{naplo}
	\naplotag{feb. 18. 16h }{ 1 óra }{ Csapat }{ Értekezlet  \newline Döntés: Segédeszközök kiválasztása (git, trello, drive)}
\end{naplo}

\chapter{Analízis modell kidolgozása 1.}

\section{Objektum katalógus}

\comment{Minden, a feladatban szereplő objektum rövid, egy-két bekezdés hosszú ismertetése. Meg kell jelenjen minden objektumhoz, hogy mi a felelőssége. Informális leírás, ezért nem kell foglalkozni az örökléssel, az interfészekkel, az absztrakt osztályokkal, a segédosztályokkal.}

\subsection{Objektum1}
\comment{Felelősség informális leírása}

\subsection{Objektum2}
\comment{Felelősség informális leírása}

\comment{Az objektumkatalógus alapján kiindulva kell megalkotni az objektumorientált analízis modellt. A 3.2, 3.3, 3.4, 3.5 alfejezetek ugyanannak a modellnek a különböző nézetei, ezért egyidőben, egymással összefüggésben készülnek.Megtörténik az objektumkatalógusbantárgyalt objektumok felelősségének formalizálásaosztályokká,attribútumokká, metódusokká. Csak publikus metódusok szerepelhetnek. Megjelennek az interfészek, az öröklés, az absztrakt osztályok. Segédosztályokra még nincs szükség.}

\section{Statikus struktúra diagramok}
\comment{Az előző alfejezet osztályainak kapcsolatait és publikus metódusait bemutató osztálydiagram(ok). Tipikus hibalehetőségek: csillag-topológia, szigetek.}

\diagram{docs/img/BMElogo}{Demó}{3cm}

\section{Osztályok leírása}
\comment{Az előző alfejezetben tárgyalt objektumok felelősségének formalizálása attribútumokká, metódusokká. Csak publikus metódusok szerepelhetnek. Ebben az alfejezetben megjelennek az interfészek, az öröklés, az absztrakt osztályok. Segédosztályokra még mindig nincs szükség. Az osztályok ABC sorrendben kövessék egymást. Interfészek esetén az Interfészek, Attribútumok pontok kimaradnak.}

\subsection{Osztály1}
\begin{class-template-responsibility}
    Felelősség leírása
\end{class-template-responsibility}
\begin{class-template-interface}
    Megvalósított interfészek felsorolása
\end{class-template-interface}
\begin{class-template-baseclass}
    Ős-Ősosztály \baseclass Ősosztály... 
\end{class-template-baseclass}
\begin{class-template-attribute}
    \classitem{+A [0..*]}{adattag A}
\end{class-template-attribute}
\comment{Milyen publikus, protected és privát  metódusokkal rendelkezik. Metódusonként egy-három mondat arról, hogy mit csinál.}
\begin{class-template-method}
    \classitem{+B(A a) : void}{metódus B}
\end{class-template-method}

\subsection{Osztály2}
\begin{class-template-responsibility}
    Felelősség leírása
\end{class-template-responsibility}
\begin{class-template-statechart}
    A belső működéshez tartozó állapotgép
    \diagram{docs/img/BMElogo}{Az osztály állapotváltozásai}{3cm}
\end{class-template-statechart}

\section{Statikus struktúra diagramok}
\comment{Az előző alfejezet osztályainak kapcsolatait és publikus metódusait bemutató osztálydiagram(ok). Tipikus hibalehetőségek: csillag-topológia, szigetek.}
\diagram{docs/img/BMElogo}{osztálydiagram}{3cm}


\section{Szekvencia diagramok}
\comment{Inicializálásra, use-case-ekre, belső működésre. Konzisztens kell legyen az előző alfejezettel. Minden metódus, ami ott szerepel, fel kell tűnjön valamelyik szekvenciában. Minden metódusnak, ami szekvenciában szerepel, szereplnie kell a valamelyik osztálydiagramon.}
\diagram{docs/img/BMElogo}{Szekvencia1}{3cm}

\section{State-chartok}
\comment{Csak azokhoz az osztályokhoz, ahol van értelme. Egyetlen állapotból álló state-chartok ne szerepeljenek. A játék működését bemutató state-chart-ot készíteni tilos.}

\clearpage
\section{Napló}

\comment{ 
    A napló tartalmazza az előző beadás óta eltelt időszak történéseit időrendben. A naplóból egyértelműen ki kell derülnie, hogy az egyes anyagrészeket ki és mennyi idő alatt készítette.
    
    A napló bejegyzésekből áll. Minden bejegyzésnek tartalmaznia kell: 
    \begin{itemize}
        \item a történés kezdetének időpontját, nap-óra pontossággal
        \item történés időtartamát, óra felbontással
        \item a szereplő(k) nevét (Kérjük a szereplők VEZETÉKNEVÉT használni)
        \item a tevékenység leírását.
    \end{itemize}
    
    Amennyiben a tevékenységben több szereplő vesz részt, akkor aza tevékenység csak értekezlet lehet, amelynek az eredményei DÖNTÉSEK. A döntéseket precízen meg kell szövegezni (Pl.: Az X objektum Y és Z metódusainak kódját W készíti el Q határidőre).Ha a bejegyzés egyetlen személyhez kötődik, akkor meg kell adni, hogy a tevékenység milyen dologra irányul. A dolog a feladat kapcsán elkészítendő termék, amelynek a (esetleg korábban) beadott anyagban megtalálhatónak kell lenni.A naplóbejegyzés felbontásának egysége szöveges, rajzos anyag esetében az ábra, diagram, vagy kb. fél-egy oldalnyi szöveg. Kódban az egység a metódus. (Pl.: A 3. ábrán látható szekvencia-diagram kidolgozása, vagy az X objektum Y és Z metódusainak kódolása és belövése.
}

\begin{naplo}
    \naplotag{feb. 18. 16h }{ 1 óra }{ Csapat }{ Értekezlet  \newline Döntés: Segédeszközök kiválasztása (git, trello, drive)}
\end{naplo}

\chapter{Analízis modell kidolgozása 2.}

\section{Objektum katalógus}

\comment{Minden, a feladatban szereplő objektum rövid, egy-két bekezdés hosszú ismertetése. Meg kell jelenjen minden objektumhoz, hogy mi a felelőssége. Informális leírás, ezért nem kell foglalkozni az örökléssel, az interfészekkel, az absztrakt osztályokkal, a segédosztályokkal.}

\subsection{Objektum1}
\comment{Felelősség informális leírása}

\subsection{Objektum2}
\comment{Felelősség informális leírása}

\comment{Az objektumkatalógus alapján kiindulva kell megalkotni az objektumorientált analízis modellt. A 3.2, 3.3, 3.4, 3.5 alfejezetek ugyanannak a modellnek a különböző nézetei, ezért egyidőben, egymással összefüggésben készülnek.Megtörténik az objektumkatalógusbantárgyalt objektumok felelősségének formalizálásaosztályokká,attribútumokká, metódusokká. Csak publikus metódusok szerepelhetnek. Megjelennek az interfészek, az öröklés, az absztrakt osztályok. Segédosztályokra még nincs szükség.}

\section{Statikus struktúra diagramok}
\comment{Az előző alfejezet osztályainak kapcsolatait és publikus metódusait bemutató osztálydiagram(ok). Tipikus hibalehetőségek: csillag-topológia, szigetek.}

\diagram{docs/img/BMElogo}{Demó}{3cm}

\section{Osztályok leírása}
\comment{Az előző alfejezetben tárgyalt objektumok felelősségének formalizálása attribútumokká, metódusokká. Csak publikus metódusok szerepelhetnek. Ebben az alfejezetben megjelennek az interfészek, az öröklés, az absztrakt osztályok. Segédosztályokra még mindig nincs szükség. Az osztályok ABC sorrendben kövessék egymást. Interfészek esetén az Interfészek, Attribútumok pontok kimaradnak.}

\subsection{Osztály1}
\begin{class-template-responsibility}
    Felelősség leírása
\end{class-template-responsibility}
\begin{class-template-interface}
    Megvalósított interfészek felsorolása
\end{class-template-interface}
\begin{class-template-baseclass}
    Ős-Ősosztály \baseclass Ősosztály... 
\end{class-template-baseclass}
\begin{class-template-attribute}
    \classitem{+A [0..*]}{adattag A}
\end{class-template-attribute}
\comment{Milyen publikus, protected és privát  metódusokkal rendelkezik. Metódusonként egy-három mondat arról, hogy mit csinál.}
\begin{class-template-method}
    \classitem{+B(A a) : void}{metódus B}
\end{class-template-method}

\subsection{Osztály2}
\begin{class-template-responsibility}
    Felelősség leírása
\end{class-template-responsibility}
\begin{class-template-statechart}
    A belső működéshez tartozó állapotgép
    \diagram{docs/img/BMElogo}{Az osztály állapotváltozásai}{3cm}
\end{class-template-statechart}

\section{Statikus struktúra diagramok}
\comment{Az előző alfejezet osztályainak kapcsolatait és publikus metódusait bemutató osztálydiagram(ok). Tipikus hibalehetőségek: csillag-topológia, szigetek.}
\diagram{docs/img/BMElogo}{osztálydiagram}{3cm}


\section{Szekvencia diagramok}
\comment{Inicializálásra, use-case-ekre, belső működésre. Konzisztens kell legyen az előző alfejezettel. Minden metódus, ami ott szerepel, fel kell tűnjön valamelyik szekvenciában. Minden metódusnak, ami szekvenciában szerepel, szereplnie kell a valamelyik osztálydiagramon.}
\diagram{docs/img/BMElogo}{Szekvencia1}{3cm}

\section{State-chartok}
\comment{Csak azokhoz az osztályokhoz, ahol van értelme. Egyetlen állapotból álló state-chartok ne szerepeljenek. A játék működését bemutató state-chart-ot készíteni tilos.}

\clearpage
\section{Napló}

\comment{ 
    A napló tartalmazza az előző beadás óta eltelt időszak történéseit időrendben. A naplóból egyértelműen ki kell derülnie, hogy az egyes anyagrészeket ki és mennyi idő alatt készítette.
    
    A napló bejegyzésekből áll. Minden bejegyzésnek tartalmaznia kell: 
    \begin{itemize}
        \item a történés kezdetének időpontját, nap-óra pontossággal
        \item történés időtartamát, óra felbontással
        \item a szereplő(k) nevét (Kérjük a szereplők VEZETÉKNEVÉT használni)
        \item a tevékenység leírását.
    \end{itemize}
    
    Amennyiben a tevékenységben több szereplő vesz részt, akkor aza tevékenység csak értekezlet lehet, amelynek az eredményei DÖNTÉSEK. A döntéseket precízen meg kell szövegezni (Pl.: Az X objektum Y és Z metódusainak kódját W készíti el Q határidőre).Ha a bejegyzés egyetlen személyhez kötődik, akkor meg kell adni, hogy a tevékenység milyen dologra irányul. A dolog a feladat kapcsán elkészítendő termék, amelynek a (esetleg korábban) beadott anyagban megtalálhatónak kell lenni.A naplóbejegyzés felbontásának egysége szöveges, rajzos anyag esetében az ábra, diagram, vagy kb. fél-egy oldalnyi szöveg. Kódban az egység a metódus. (Pl.: A 3. ábrán látható szekvencia-diagram kidolgozása, vagy az X objektum Y és Z metódusainak kódolása és belövése.
}

\begin{naplo}
    \naplotag{feb. 18. 16h }{ 1 óra }{ Csapat }{ Értekezlet  \newline Döntés: Segédeszközök kiválasztása (git, trello, drive)}
\end{naplo}

\chapter{Szkeleton tervezése}


\section{A szkeleton modell valóságos use-case-ei}
\comment{A szkeletonnak, mint önálló programnak a működésével kapcsolatos use-case-ek.}

\subsection{Use-case diagram}
\diagram{docs/img/BMElogo}{Demó}{4cm}

\subsection{Use-case leírások}
\comment{Minden use-case-hez külön}
\begin{use-case}
    %név
    {Use-case Neve}
    %rövid leírás
    {Az eset rövid leírása}
    %aktorok
    {Aktorok}
    %forgatókönyv
    Forgatókönyv \newline 
    \textbf{A.1} Alternatíva
\end{use-case}

\section{A szkeleton kezelői felületének terve, dialógusok}
\comment{A szkeleton által elfogadott bemenetek , valamint a szöveges konzolon megjelenő kimenetek. A kiemenet formátuma olyan kell legyen, ami alapján a működés összevethető a korábbi szekvencia-diagramokkal.}

\comment{Konzol input/output szemléltetésre:}
\begin{verbatim}
    move 1 to 3
    get item
\end{verbatim}

\section{Szekvencia diagramok a belső működésre}
\comment{A szkeletonban implementált szekvenciadiagramok. Tipikusan egy use-case egy diagram. Ezek megegyezhetnek a korábban specifikált diagramokkal, de az egyes életvonalakat (lifeline) egyértelműen a szkeletonban példányosított objektumokhoz kell tudni kötni. Azt kell megjeleníteni, hogy a szkeletonban létrehozott objektumok egymással hogyan fognak kommunikálni.}

\section{Kommunikációs diagramok}
\comment{A szkeletonban, az egyes szkeleton-use-case-ek futása során létrehozott objektumok és kapcsolataik bemutatására szolgáló diagramok. Ezek alapján valósítják meg a szkeleton fejlesztői az inicializáló kódrészleteket.}

\clearpage
\section{Napló}

\comment{ 
    A napló tartalmazza az előző beadás óta eltelt időszak történéseit időrendben. A naplóból egyértelműen ki kell derülnie, hogy az egyes anyagrészeket ki és mennyi idő alatt készítette.
    
    A napló bejegyzésekből áll. Minden bejegyzésnek tartalmaznia kell: 
    \begin{itemize}
        \item a történés kezdetének időpontját, nap-óra pontossággal
        \item történés időtartamát, óra felbontással
        \item a szereplő(k) nevét (Kérjük a szereplők VEZETÉKNEVÉT használni)
        \item a tevékenység leírását.
    \end{itemize}
    
    Amennyiben a tevékenységben több szereplő vesz részt, akkor aza tevékenység csak értekezlet lehet, amelynek az eredményei DÖNTÉSEK. A döntéseket precízen meg kell szövegezni (Pl.: Az X objektum Y és Z metódusainak kódját W készíti el Q határidőre).Ha a bejegyzés egyetlen személyhez kötődik, akkor meg kell adni, hogy a tevékenység milyen dologra irányul. A dolog a feladat kapcsán elkészítendő termék, amelynek a (esetleg korábban) beadott anyagban megtalálhatónak kell lenni.A naplóbejegyzés felbontásának egysége szöveges, rajzos anyag esetében az ábra, diagram, vagy kb. fél-egy oldalnyi szöveg. Kódban az egység a metódus. (Pl.: A 3. ábrán látható szekvencia-diagram kidolgozása, vagy az X objektum Y és Z metódusainak kódolása és belövése.
}

\begin{naplo}
    \naplotag{feb. 18. 16h }{ 1 óra }{ Csapat }{ Értekezlet  \newline Döntés: Segédeszközök kiválasztása (git, trello, drive)}
\end{naplo}

\documentclass[../../projlab]{subfiles}
\begin{document}

\makeatletter

\ifSubfilesClassLoaded{
	\coverpage{5. Leadás}
	\def\filePath[#1]{./../../#1}
}{}

\makeatother

\chapter{Szkeleton beadás}

\section{Fordítási és futtatási útmutató}
\comment{A feltöltött program fordításával és futtatásával kapcsolatos útmutatás. Ennek tartalmaznia kell leltárszerűen az egyes fájlok pontos nevét, méretét byte-ban, keletkezési idejét, valamint azt, hogy a fájlban mi került megvalósításra.}

\subsection{Fájllista}

\begin{fajllista}
	\fajl
	{Demo.java} % Kezdet
	{353 byte} % Méret
	{2020.03.26~21:05~} % Keletkezés ideje
	{Demo programosztály} % Leírás
\end{fajllista}

\subsection{Fordítás}
\comment{A fenti listában szereplő forrásfájlokból milyen műveletekkel lehet a bináris, futtatható kódot előállítani. Az előállításhoz csak a 2. Követelmények c. dokumentumban leírt környezetet szabad előírni.}

\begin{verbatim}
    javac -d bin *.java
\end{verbatim}

\subsection{Futtatás}
\comment{A futtatható kód elindításával kapcsolatos teendők leírása. Az indításhoz csak a 2. Követelmények c. dokumentumban leírt környezetet szabad előírni.}

\begin{verbatim}
    cd bin
    java Main.java
\end{verbatim}

\clearpage
\section{Értékelés}
\comment{A projekt kezdete óta az értékelésig eltelt időben tagokra bontva, százalékban.}
\begin{ertekeles}
	  \ertekelestag{Teszt János}{ ?????? }{100\%}
\end{ertekeles}


\end{document}
\clearpage
\section{Napló}

\comment{ 
    A napló tartalmazza az előző beadás óta eltelt időszak történéseit időrendben. A naplóból egyértelműen ki kell derülnie, hogy az egyes anyagrészeket ki és mennyi idő alatt készítette.
    
    A napló bejegyzésekből áll. Minden bejegyzésnek tartalmaznia kell: 
    \begin{itemize}
        \item a történés kezdetének időpontját, nap-óra pontossággal
        \item történés időtartamát, óra felbontással
        \item a szereplő(k) nevét (Kérjük a szereplők VEZETÉKNEVÉT használni)
        \item a tevékenység leírását.
    \end{itemize}
    
    Amennyiben a tevékenységben több szereplő vesz részt, akkor aza tevékenység csak értekezlet lehet, amelynek az eredményei DÖNTÉSEK. A döntéseket precízen meg kell szövegezni (Pl.: Az X objektum Y és Z metódusainak kódját W készíti el Q határidőre).Ha a bejegyzés egyetlen személyhez kötődik, akkor meg kell adni, hogy a tevékenység milyen dologra irányul. A dolog a feladat kapcsán elkészítendő termék, amelynek a (esetleg korábban) beadott anyagban megtalálhatónak kell lenni.A naplóbejegyzés felbontásának egysége szöveges, rajzos anyag esetében az ábra, diagram, vagy kb. fél-egy oldalnyi szöveg. Kódban az egység a metódus. (Pl.: A 3. ábrán látható szekvencia-diagram kidolgozása, vagy az X objektum Y és Z metódusainak kódolása és belövése.
}

\begin{naplo}
    \naplotag{feb. 18. 16h }{ 1 óra }{ Csapat }{ Értekezlet  \newline Döntés: Segédeszközök kiválasztása (git, trello, drive)}
\end{naplo}

\documentclass[../../projlab]{subfiles}
\begin{document}

\makeatletter

\ifSubfilesClassLoaded{
	\coverpage{7. Leadás}
	%\renewcommand{\filePath}[1]{./../../ #1}
	\def\filePath[#1]{./../../#1}
}{}

\makeatother


\chapter{Prototípus koncepciója}
\setcounter{section}{-1}

\section{Változás hatása a modellre}

\subsection{Módosult osztálydiagram}
\begin{figure}[H] 
    \centering 
    \includegraphics[width=1\textwidth]{docs/bead07/Diagrams/class-diag.eps} 
    \caption{} 
\end{figure} 

\subsection{Új vagy megváltozott metódusok}

\textbf{A bevezetett új osztályok}

\subsubsection{IInvertory interface}
\begin{class-template-responsibility}
    A különböző tárolók közös interface-e
\end{class-template-responsibility}
\begin{class-template-method}
    \classitem{+getItems(): List<Item>}
                {Visszaadja a tárolt nyersanyagok listáját.}
    \classitem{+tryInsertItem(i: Item): bool}
                {Visszaadja, hogy az adott tárgy belehelyezhető-e a tárhelybe.}
    \classitem{+insertItem(i: Item): bool}
                {Az adott tárgyat belehelyezzük a tárhelybe.}
    \classitem{+removeItem(i: Item): bool}
                {Az adott tárgyat kiveszzük.}                
\end{class-template-method}

\subsubsection{AsteroidInvertory}
\begin{class-template-responsibility}
    Az aszteroida tárhelyének reprezentációja.
\end{class-template-responsibility}
\begin{class-template-interface}
    IInventory
\end{class-template-interface}
\begin{class-template-attribute}
    \classitem{+items: List<Item>}{A tárolt nyersanyagok listája}
\end{class-template-attribute}

\subsubsection{SSInvertory}
\begin{class-template-responsibility}
    Az űrhajók tárhelyének reprezentációja.
\end{class-template-responsibility}
\begin{class-template-interface}
    IInventory
\end{class-template-interface}
\begin{class-template-attribute}
    \classitem{+items: List<Item>}{A tárolt nyersanyagok listája}
    \classitem{+tgs: List<TeleportGateItem>}{A tárolt kapuk listája}
\end{class-template-attribute}

\subsubsection{InfiniteInvertory}
\begin{class-template-responsibility}
    Az ufo tárhelyének reprezentációja. Bármennyi tárgy eltárolható benne.
\end{class-template-responsibility}
\begin{class-template-interface}
    IInventory
\end{class-template-interface}
\begin{class-template-attribute}
    \classitem{+items: List<Item>}{A tárolt nyersanyagok listája}
\end{class-template-attribute}

\subsubsection{UFO}
\begin{class-template-responsibility}
    Az ufokat reprezentálja.
    Az aszteroidaöv nyersanyagaira fáj a foga, azokat akarja eltulajdonítani, de nem tud fúrni. Éppen ezért a telepesek ellenségeinek tekinthető.
\end{class-template-responsibility}
\begin{class-template-baseclass}
    Entity \baseclass MovableEntity \baseclass Vessel
\end{class-template-baseclass}
\begin{class-template-attribute}
    \classitem{\#getInvertory(): IInvertory}{Visszaadja az UFO raktárát.}
\end{class-template-attribute}

\textbf{A korábbi osztályokban történt módosítások}

\subsubsection{SpaceShip}
\begin{class-template-responsibility}
    Egy telepest reprezentáló osztály.
\end{class-template-responsibility}
\begin{class-template-baseclass}
    Entity \baseclass MovableEntity \baseclass Vessel
\end{class-template-baseclass}
\begin{class-template-attribute}
    \classitem{\#lastPlacedGate}{Törölve}
\end{class-template-attribute}
\begin{class-template-method}
    \classitem{\#getInvertory(): IInvertory}{Visszaadja a telepes által folyamatosan hordozott raktárat.}
\end{class-template-method}


\subsubsection{TeleportGate}
\begin{class-template-responsibility}
    A teleportkapukat reprezentáló osztály.
\end{class-template-responsibility}
\begin{class-template-baseclass}
    Entity \baseclass MovableEntity \baseclass Building
\end{class-template-baseclass}
\begin{class-template-attribute}
    \classitem{\-gateID: int}{A kapuk összekötéséhez használt ID}
\end{class-template-attribute}

\subsubsection{Building}
\begin{class-template-responsibility}
    Az épülettípusok közös ősosztálya.
\end{class-template-responsibility}
\begin{class-template-baseclass}
    Entity \baseclass MovableEntity  (Új ősosztály)
\end{class-template-baseclass}

\subsubsection{Entity}
\begin{class-template-responsibility}
    Egy adott entitás ( vizuális megjelenítéssel rendelkező játékelem) osztálya. 
\end{class-template-responsibility}
\begin{class-template-method}
    \classitem{+explode(): void}{Az adott entitás felrobban.}
\end{class-template-method}

\subsubsection{AsteroidPlaces (enum)}
\begin{class-template-responsibility}
    Régi név: BuildingPlace
    Azt írja le hogy egy MovableEntity hol helyezkedhet el az aszteroidán.
\end{class-template-responsibility}
\begin{class-template-attribute}
    \classitem{inside}{Az aszteroida belseje }
    \classitem{orbit}{Az aszteroida körül fix pályán}
    \classitem{vessel}{Az aszteroida mellet lévő járművek}
\end{class-template-attribute}

\subsubsection{Asteroid}
\begin{class-template-responsibility}
    Egy aszteroidát jelöl. Minden példányának létrejöttekor beállítódnak a generátor által a tulajdonságai.
\end{class-template-responsibility}
\begin{class-template-attribute}
    \classitem{stationed: Vessel[0..*]}{A jelenleg itt lévő Vesselek}
    \classitem{orbit: MovableEntity[0..*]}{A jelenleg itt pályán lévő dolgok}
    \classitem{inside: MovableEntity[0..*]}{A jelenleg ebben az aszetoidában lévő dolgok}
\end{class-template-attribute}


\subsubsection{Vessel}
\begin{class-template-responsibility}
    Egy-egy járműért (pl. telepes vagy robot) felelős osztály. 
    A bányászást is ebbe az osztályba mozgattuk.
\end{class-template-responsibility}
\begin{class-template-method}
    \classitem{\#getInvertory(): IInvertory}{Visszaadja az űrjármű által folyamatosan hordozott raktárat.}
    \classitem{+mine(): void}{A függvény megegyezik a régen a SpaceShip osztályban lévővel, az aszteroidából fogunk tudni nyersanyagot kivenni.}
\end{class-template-method}


\subsubsection{Invertory}
\begin{class-template-responsibility}
    Törölve
\end{class-template-responsibility}

\subsubsection{Item}
\begin{class-template-attribute}
    \classitem{amount: int}{Törölve}
    \classitem{inv: IInvnetory}{A tároló amiben jelenleg vagyunk}
\end{class-template-attribute}
\begin{class-template-method}
    \classitem{+nearSun(): void}{Mi történjen, ha közel kerül a nyersanyag a naphoz, és napfény éri.}
\end{class-template-method}

\subsubsection{Resource}
\begin{class-template-baseclass}
    Item
\end{class-template-baseclass}
\begin{class-template-attribute}
    \classitem{amount: int}{Törölve}
\end{class-template-attribute}

\subsubsection{Iron}
\begin{class-template-baseclass}
    Item
\end{class-template-baseclass}

\subsubsection{Ice}
\begin{class-template-baseclass}
    Item
\end{class-template-baseclass}

\subsubsection{Coal}
\begin{class-template-baseclass}
    Item
\end{class-template-baseclass}

\subsubsection{Titanium}
\begin{class-template-responsibility}
    Régi név: Titan
\end{class-template-responsibility}
\begin{class-template-baseclass}
    Item
\end{class-template-baseclass}

\subsubsection{Uranium}
\begin{class-template-baseclass}
    Item
\end{class-template-baseclass}
\begin{class-template-attribute}
    \classitem{+counter: int}{Számolja, hogy hányszor lett kitéve napsugárzásnak az adott nyersanyag.}
\end{class-template-attribute}

\subsubsection{IronResorce}
\begin{class-template-responsibility}
    Törölve
\end{class-template-responsibility}

\subsubsection{IceResource}
\begin{class-template-responsibility}
    Törölve
\end{class-template-responsibility}

\subsubsection{CoalResource}
\begin{class-template-responsibility}
    Törölve
\end{class-template-responsibility}

\subsubsection{TitaniumResorce}
\begin{class-template-responsibility}
    Törölve
\end{class-template-responsibility}

\subsubsection{UraniumResource}
\begin{class-template-responsibility}
    Törölve
\end{class-template-responsibility}

\subsubsection{TeleportGateRecipe}
\begin{class-template-attribute}
    \classitem{amount: int}{Törölve}
\end{class-template-attribute}
\begin{class-template-method}
    \classitem{+CanCraft(Invertory inv, Asteroid a): bool}
                {Az ősosztály metódusát írjuk felül, mert ellenőrizni kell a telepes tárhelyében kévő teleportkapujainak számát.}
\end{class-template-method}

\subsection{Szekvencia-diagramok}

%\comment{Az analízis modell szekvenciadiagramjaiból a változás által érintett, előírt, módosítottdiagramok}

\begin{figure}[H] 
    \centering 
    \includegraphics[width=0.7\textwidth]{docs/bead07/Diagrams/TG-SolarFlare.eps} 
    \caption{} 
\end{figure}

\begin{figure}[H] 
    \centering 
    \includegraphics[width=0.7\textwidth]{docs/bead07/Diagrams/Uranium/uran-counter} 
    \caption{} 
\end{figure}

\begin{figure}[H] 
    \centering 
    \includegraphics[width=0.7\textwidth]{docs/bead07/Diagrams/SolarFlare/solar-flare} 
    \caption{} 
\end{figure}

\begin{figure}[H] 
    \centering 
    \includegraphics[width=1\textwidth]{docs/bead07/Diagrams/Craft/CanCraft} 
    \caption{A recepthez szükséges nyersanyagokat egyenként, és nem típusonként tároljuk, ezért megváltoztak az ellenőrzés lépései. A két raktár nyersanyagkészletét lemásoljuk, és ha megtaláljuk a keresett összetevőket, eltávolítjuk a lemásolt listákból, így az eredeti lista tartalma sikertelen készítés esetén sem fog változni. } 
\end{figure} 

\begin{figure}[H] 
    \centering 
    \includegraphics[width=1\textwidth]{docs/bead07/Diagrams/Craft/CanCraftGate} 
    \caption{A teleportkapuk előállításának ellenőrzése speciális eset lesz, mert itt ellenőrizni kell, hogy elférnek-e a telepes raktárában a kapuk készítés után.} 
\end{figure} 

\begin{figure}[H] 
    \centering 
    \includegraphics[width=1\textwidth]{docs/bead07/Diagrams/Craft/Craft} 
    \caption{A raktárak megváltozott elemkészlete miatt (minden nyersanyag külön van tárolva) az ellenőrzés, eltávolítás folyamata megváltozott, egyenként távolítjuk el az összetevőket. } 
\end{figure} 

\begin{figure}[H] 
    \centering 
    \includegraphics[width=0.7\textwidth]{docs/bead07/Diagrams/Craft/Gate_MakeResult} 
    \caption{A teleportkapu-párok tagjait külön tároljuk, így 2 kaput hozunk létre, és hozzájuk rendelünk egy közös azonosítót.} 
\end{figure} 

\begin{figure}[H] 
    \centering 
    \includegraphics[width=0.5\textwidth]{docs/bead07/Diagrams/Inventory/InsertItemtoInfinite} 
    \caption{Egy tárgy elhelyezése az UFO tárolójába} 
\end{figure} 

\begin{figure}[H] 
    \centering 
    \includegraphics[width=0.5\textwidth]{docs/bead07/Diagrams/Inventory/InsertItemtoSSinventory} 
    \caption{Egy tárgy elhelyezése a SpaceShip tárolójába} 
\end{figure} 

\begin{figure}[H] 
    \centering 
    \includegraphics[width=0.7\textwidth]{docs/bead07/Diagrams/Inventory/InsertTeleportGatetoInventory} 
    \caption{Egy TeleportGate elhelyezése az SpaceShip tárolójába} 
\end{figure}

\begin{figure}[H] 
    \centering 
    \includegraphics[width=0.7\textwidth]{docs/bead07/Diagrams/Mine/Vessel-mine} 
    \caption{A játékos bányászik egy aszteroidából} 
\end{figure} 

\begin{figure}[H] 
    \centering 
    \includegraphics[width=0.7\textwidth]{docs/bead07/Diagrams/Move/Move-ME-to-asteroid} 
    \caption{A MoveableEntity mozgatása} 
\end{figure} 

\begin{figure}[H] 
    \centering 
    \includegraphics[width=0.7\textwidth]{docs/bead07/Diagrams/Explode/Uran-explode} 
    \caption{A MoveableEntity mozgatása} 
\end{figure} 


\section{Prototípus interface-definíciója}

Az alkalmazás minden parancsot a standard input-ról olvas be és értelmez. A teszeseteket a meghatározott módon lehet majd végrehajtani.
A kimenet a standard output-ra kerül kiíratásra.
\subsection{Az interfész általános leírása}

A program standard input-ot használva vezérelhető a megadott parancsok segítségével. Minden tesztesetet a tesztelő a meghatározott parancsok és paraméterek előre definiált sorrendben való kiadásával hajtja végre. Egyedi tesztesetek is létrehozhatóak, ezek ellenőrzésért a tesztelő felelős.
Minden esetben a kimenet standard output-ra kerül kiíratásra.
\subsection{Bemeneti nyelv}


\subsubsection{Bemeneti parancsok}

\begin{itemize}
    \item startgame
    \begin{itemize}
        \item Leírás: 
        \item Opciók: startgame <seed> <input | generate>
        \begin{itemize}
            \item seed: a véletlen szám generátor által hasznát kezdőérték.
            \item input: A felhasználó adja meg a játéktér összeállítását
            \item generate: A játéktér véletlenszerűen generálódik.
        \end{itemize}
    \end{itemize}

    \item move
    \begin{itemize}
        \item Leírás: A játékos egy járművet akar mozgatni. \newline
            A parancs megkérdezi a játékostól, hogy melyik aszteroidára akarja mozgatni a járművet.
        \item Opciók: move <vessel>
        \begin{itemize}
            \item vessel: A jármű amit mozgatni akarunk
        \end{itemize}
        \item Példa:
            \begin{verbatim}
> move ss_001
Neighbours:
01 - Asteroid_1  - ?
02 - Asteroid_34 - Iron
03 - Asteroid_9  - Coal
04 - Asteroid_11 - ?
> 04
ss_001 Moved to Asteroid_11
            \end{verbatim}
    \end{itemize}


    \item neighbours
    \begin{itemize}
        \item Leírás: A játékos lekéri az aszteroida szomszédjait
        \item Opciók: neighbours <asteroid>
        \begin{itemize}
            \item asteroid: Az aszteroida aminek a szomszédjait akarjuk lekérdezni.
        \end{itemize}
        \item Példa:
            \begin{verbatim}
> neighbours ss_001
Neighbours:
01 - Asteroid_1  - ?
02 - Asteroid_34 - Iron
03 - Asteroid_9  - Coal
04 - Asteroid_11 - ?
            \end{verbatim}
    \end{itemize}


    \item drill
    \begin{itemize}
        \item Leírás: A játékos a megadott járművel fúr azon az aszteroidán amin éppen van.
        \item Opciók: drill <vessel>
        \begin{itemize}
            \item vessel: A jármű amivel fúrni akarunk.
        \end{itemize}
        \item Példa 1:
            \begin{verbatim}
> drill ss_001 
Asteroid_42 crust = X
            \end{verbatim}
        \item Példa 2:
            \begin{verbatim}
> drill ss_001 
can't drill
            \end{verbatim}
    \end{itemize}


    \item mine
    \begin{itemize}
        \item Leírás: A játékos a megadott járművel megpróbál bányászni azon az aszteroidán amin éppen van.
        \item Opciók: mine <vessel>
        \begin{itemize}
            \item vessel: A jármű amivel bányászni akarunk.
        \end{itemize}
        \item Példa:
            \begin{verbatim}
> mine ss_001
Asteroid_35 mined
* Item: Iron
            \end{verbatim}
    \end{itemize}


    \item status
    \begin{itemize}
        \item Leírás: A játékos a megadott jármű jelenlegi állapotát kérdezi le.
        \item Opciók: status <vessel|asteroid>
        \begin{itemize}
            \item vessel: A jármű aminek az állapotát kérdezzük le.
            \item asteroid: Az aszteroida aminek az állapotát kérdezzük le.
        \end{itemize}
        \item Példa 1:
            \begin{verbatim}
> status ss_001
Asteroid: Asteroid_67 | Not hidden
Resources:
- Iron
- Iron
- Coal
- etc....
Teleport Gates:
- Tg_023 - 1
- Tg_023 - 1
- Tg_011 - 5
            \end{verbatim}
        \item Példa 2:
            \begin{verbatim}
> status Asteroid_67
Crust: 0
Vessels:
- SpaceShip_01
- SpaceShip_23
- Robot_4
Resources:
- Iron
- Iron
Teleport Gates:
- Tg_023 - 1
- Tg_011 - 5
            \end{verbatim}        
    \end{itemize}


    \item hide
    \begin{itemize}
        \item Leírás: A játékos 
        \item Opciók: hide <enter|exit> <vessel>
        \begin{itemize}
            \item vessel: A jármű aminek az állapotát kérdezzük le.
        \end{itemize}
        \item Példa:
            \begin{verbatim}
> hide enter ss_001
Can't hide
            \end{verbatim}
        \item Példa:
            \begin{verbatim}
> hide enter ss_001
ss_001 is hidden
            \end{verbatim}            
    \end{itemize}


    \item craft
    \begin{itemize}
        \item Leírás: A játékos egy űrhajón legyárt egy adott receptet.
        \item Opciók: craft <vessel> <recipe>
        \begin{itemize}
            \item vessel: A jármű amin akarunk gyártani
            \item recipe: A recept amit le akarunk gyártani
        \end{itemize}
        \item Példa:
            \begin{verbatim}
> craft ss_001 <recipe>
Can't craft .......
--
Succesful
            \end{verbatim}
    \end{itemize}


    \item place
    \begin{itemize}
        \item Leírás: A játékos egy űrhajóból lerak egy nyersanyagot az aszteroidába amin van.
        \item Opciók: place <resource> from <vessel>
        \begin{itemize}
            \item resource: a nyersanyag amit le akarunk rakni
            \item vessel: A jármű amin akarunk gyártani
        \end{itemize}
        \item Példa:
            \begin{verbatim}
> palce Iron from ss_001
            \end{verbatim}
    \end{itemize}

    \item pickup
    \begin{itemize}
        \item Leírás: A játékos egy űrhajóból lerak egy nyersanyagot az aszteroidába amin van.
        \item Opciók: pickup <resource> to <vessel>
        \begin{itemize}
            \item resource: a nyersanyag amit le akarunk rakni
            \item vessel: A jármű amin akarunk gyártani
        \end{itemize}
        \item Példa:
            \begin{verbatim}
> pickup Coal to ss_001
            \end{verbatim}
    \end{itemize}
\end{itemize}

\subsubsection{Események}
Továbbá a program bizonyos események után lehet hogy eseményekről értesít minket.
Például:
\begin{verbatim}
[Event] Asteroid_42 exploded
[Event] ss_001 exploded
\end{verbatim}

%\comment{Ha szükséges, meg kell adni a konfigurációs (pl. pályaképet megadó) fájlok nyelvtanát is.}

\subsection{Kimeneti nyelv}
%\comment{Egyértelműen definiálni kell, hogy az egyes bemeneti parancsok végrehajtása után előálló állapot milyen formában jelenik meg a szabványos kimeneten. A program képes legyen olyan kimenetet előállítani, amellyel az objektumok állapota ellenőrizhető (pl. save). Ebben az alfejezetben is precízen definiálni kell, hogy a kimenet nyelve milyen elemekből és milyen szintakszissal áll elő.}

A kimeneten minden parancs végrehajtása után megjelenik a változtatott objektumok állapota. Ezen kívűl a futás közben is lekérdezhetőek ezek az adatok a status parancsot használva.
Az így előállt kimenet a következő mintát követi:
\begin{verbatim}
...
Osztály: Név
Belső változók: értékek
...
\end{verbatim}

Ezen kívűl az előző pontban is említett [Event]-el ellátott értesítések is kiírásra kerülnek.
pl.:
\begin{verbatim}
...
[Event] asteroid exploded
...
\end{verbatim}



\section{Összes részletes use-case}
\begin{use-case}
    {startgame}
    {A játékos elindítja a játékot.}
    {Játékos} 
    \textbf{A} A játékos kiad egy startgame <seed> input parancsot\newline
    \textbf{A.1} Az indítás sikeres \newline
    \textbf{A.2} Az indítás sikertelen \newline
    \textbf{B} A játékos kiad egy startgame <seed> generate parancsot\newline
	\textbf{B.1} Az indítás sikeres \newline
    \textbf{B.2} Az indítás sikertelen 

\end{use-case}

\begin{use-case}
    {move}
    {A játékos mozgatja az egyik járművet.}
    {Játékos} 
    \textbf{A} A játékos kiadja a move <vessel> <to> parancsot \newline
	\textbf{A.1} Az mozgatás sikeres \newline
    \textbf{A.2} Az mozgatás sikertelen

\end{use-case}

\begin{use-case}
    {neighbours}
    {A játékos megnézi egy aszteroida környezetét.}
    {Játékos} 
    \textbf{A} A játékos kiadja a neighbours <asteroid> parancsot \newline
	\textbf{A.1} Az elérhető aszteroidák kiíratásra kerülnek

\end{use-case}

\begin{use-case}
    {drill}
    {A játékos egy járművel próbál fúrni.}
    {Játékos} 
    \textbf{A} A fúrás sikeres \newline
    \textbf{B} A fúrás sikertelen (Már át van fúrva a kéreg) \newline
    \textbf{B} A fúrás sikertelen (A jármű nem tud fúrni) \newline
\end{use-case}

\begin{use-case}
    {mine}
    {A játékos egy járművel próbál bányászni.}
    {Játékos} 
    \textbf{A} A bányászás sikeres \newline
    \textbf{B.1} A bányászás sikertelen (Nincs elég tárhely) \newline
    \textbf{B.2} A bányászás sikertelen (Nincs mit bányászni) \newline
    \textbf{B.3} A bányászás sikertelen (Nincs átfúrva a kéreg) \newline
    \textbf{B.3} A bányászás sikertelen (A jármű nem tud bányászni) \newline
\end{use-case}

\begin{use-case}
    {status}
    {A játékos lekérdezi egy jármű vagy aszteroida státuszát}
    {Játékos} 
    \textbf{A} A játékos lekérdezi egy Jármű státuszát \newline
    \textbf{B} A játékos lekérdezi egy Aszteroida státuszát
\end{use-case}

\begin{use-case}
    {hide}
    {A játékos megkér egy űrhajót arra hogy elbujjon.}
    {Játékos} 
    \textbf{A} A játékos kiad egy hide parancot \newline
    \textbf{A.1} Sikeresen elbújik a jármű \newline
    \textbf{A.2} A jármű nem tud elbújni.
\end{use-case}

\begin{use-case}
    {craft}
    {A játékos legyártat egy receptaet egy űrhajón}
    {Játékos} 
    \textbf{A} A játékos kiad egy craft parancot \newline
    \textbf{A.1} A gyártás sikeres \newline
    \textbf{A.2.a} A gyártás sikertelen (nincs elég nyersanyag)
    \textbf{A.2.b} A gyártás sikertelen (nincs elég tárhely)
\end{use-case}

\begin{use-case}
    {place}
    {A játékos lerakat egy nyersanyagot az aszteroidába}
    {Játékos} 
    \textbf{A} A játékos kiad egy place parancot \newline
    \textbf{A.1} A lerakás sikeres \newline
    \textbf{A.2} A lerakás sikertelen 
\end{use-case}

\begin{use-case}
    {pickup}
    {A játékos felvetet egy nyersanyagot az aszteroidából}
    {Játékos} 
    \textbf{A} A játékos kiad egy pickup parancot \newline
    \textbf{A.1} A felvétel sikeres \newline
    \textbf{A.2} A felvétel sikertelen 
\end{use-case}


\section{Tesztelési terv}
%\comment{A tesztelési tervben definiálni kell, hogy a be- és kimeneti fájlok egybevetésével miként végezhető el a program tesztelése. Meg kell adni teszt forgatókönyveket. Az egyes teszteket elég informálisan, szabad szövegként leírni. Teszt-esetenként egy-öt mondatban. Minden teszthez meg kell adni, hogy mi a célja, a proto mely funkcionalitását, osztályait stb. teszteli. Az alábbi táblázat minden teszt-esethez külön-külön elkészítendő.}

\test{Move movable entity}
{Mozgatható entitás mozgatása egy szomszédos aszteroidára}
{Sikeresen mozgás a pályán}

\test{Vessel moves through teleport gate}
{Járművek mozgása egy szomszédos aszteroidára teleportkapun keresztül}
{Teleportkapuk szomszédossá tesznek két aszteroidát.}

\test{Robot/SpaceShip drills}
{Robot/Telepes fúr az adott aszteroida köpenyébe, mélyíti azt.}
{Robot/Telepes segít átfúrni az aszteroidákat.}

\test{Robot/SpaceShip can't drill because of core size}
{A robot/telepes próbál fúrni egy aszteroida köpenyébe, de nem lehetséges, mert át van fúrva.}
{A robot/telepes lépése meghiúsul, egy átfúrt magú aszteroida nem fúrható tovább. }

\test{Robot/SpaceShip can't drill because of explosion}
{A robot/telepes próbál fúrni egy aszteroida köpenyébe, de nem lehetséges, mert felrobbant az aszteroida.}
{A robot/telepes lépése meghiúsul, egy felrobbant aszteroida nem fúrható tovább. }

\test{Craft robot}
{A telepes a szükséges nyersanyagokkal egy robotot készít el. }
{Az adott nyersanyagokból elkészül egy robot, amit az adott játékos irányít. }

\test{Can't craft robot}{A telepes egy robotot szeretne elkészíteni, de nincs elegendő nyersanyaga hozzá. }{Robot sikertelen készítése}
\test{Craft space station}{A telepes a szükséges nyersanyagokkal egy űrállomást készít el, hogy megnyerje a játékot. }{Az adott nyersanyagokból elkészül egy űrállomás, ami a játék végét jelenti. }
\test{Can't craft space station}{A telepes egy űrállomást szeretne elkészíteni, de nincs elegendő nyersanyaga hozzá. }{Űrállomás sikertelen készítése}
\test{Craft teleport gate}{A telepes a szükséges nyersanyagokkal egy teleportkapu-párt készít el. }{Az adott nyersanyagokból elkészül egy kapupár, amit az adott játékos elhelyez a tárhelyében. }
\test{Can't craft teleport gate}{A telepes egy teleportkapu-párt szeretne elkészíteni, de nincs elegendő nyersanyaga hozzá. }{Teleportkapu-pár sikertelen készítése}
\test{Not enough space to craft teleport gate}{A telepes raktárában nincs elég hely a teleportkapu-pár tárolására.}{Sikertelen teleportkapu-pár készítés, tárhely ellenőrzés sikeressége. }
\test{Direct robots}{Robotoknak megadja a játékos, hogy milyen nyersanyag után kutassanak.} {Robotok célzott irányítása nyersanyagkutatás céljából, kért nyersanyag tárolása.}
\test{SpaceShip remains hidden}{A telepes meg szeretne bújni, amikor már meg van bújva.}{A megbújás állapota mindaddig megmarad, amíg ki nem bújik. }
\test{SpaceShip remains unhidden}{A telepes ki szeretne bújni, amikor már ki van bújva.}{A kibújás állapota mindaddig megmarad, amíg meg nem bújik. }
\test{SpaceShip can't hide in exploded asteroid}{A telepes el szeretne bújni, de az aszteroida fel van robbanva.}{Sikertelen megbújás, felrobbant aszteroida magjába nem lehet megbújni. }
\test{SpaceShip can't hide in undrilled asteroid}{A telepes el szeretne bújni, de az aszteroida magja nincs teljesen átfúrva.}{Sikertelen megbújás, nem átfúrt aszteroida magjába nem lehet megbújni.}
\test{SpaceShip doesn't have space to hide}{Egy átfúrt magú aszteroidában meg szeretne bújni a telepes, de a mag "meg van telve".}{Sikertelen megbújás, egy átfúrt aszteroida magjában csak limitált számú telepes fér el. }
\test{SpaceShip hides}{A telepes el szeretne, és el tud bújni. }{A telepes az adott körülmények ellenőrzése után sikeresen meg tud bújni egy aszteroida magjában.}
\test{SpaceShip exits hiding}{A telepes meg van bújva, és kibújik az aszteroida magjából.}{Sikeres művelet, a telepesek ki tudnak bújni az aszteroidából és folytatni tevékenységeiket.}
\test{SpaceShip mines/UFO}{A telepes/UFO egy átfúrt magú aszteroidából tud bányászni, ha van benne elérhető nyersanyag.}{Sikeres nyersanyagbányászat művelete, a nyersanyag a telepes/UFO raktárába kerül.}
\test{SpaceShip/UFO can't mine because of explosion}{Egy felrobbant aszteroida magjában próbál bányászni a telepes.}{Sikertelen művelet, felrobbant aszteroidából nem lehet bányászni.}
\test{SpaceShip/UFO can't mine because of empty core}{Egy telepes/UFO bányászni szeretne egy üres magú aszteroidából.}{Sikertelen művelet, üres magú aszteroidából nem lehet bányászni. }
\test{SpaceShip/UFO can't mine because of crust size}{Egy telepes/UFO egy még nem átfúrt magú aszteroidából próbál bányászni. }{Sikertelen művelet, nem átfúrt magú aszteroidából nem lehet bányászni.}
\test{SpaceShip can't mine because of full inventory}{A telepes ki tudja hozni a nyersanyagot a magból, de tele van a raktára.}{Sikertelen művelet, ha tele van a raktár, nem tud bányászni.}
\test{Move asteroid field}{A Naptól mért távolság változik.}{Az aszteroidák napközeli és nem napközeli állapotokban lehetnek.}
\test{Place gate}{A telepes a raktárában lévő teleportkapu-párt le szeretné helyezni.}{Egy teleportkapu-pár sikeres elhelyezése, szomszédságok beállítása, és raktárból való eltávolítása.}
\test{Place resource in asteroid}{Az aszteroida raktárába a telepes kibányászott nyersanyagot szeretne elhelyezni.}{Sikeres elhelyezés, a nyersanyag áthelyezhető a telepes raktárából az aszteroidába. }
\test{Can't place resource in exploded asteroid}{A felrobbant aszteroida raktárába a telepes kibányászott nyersanyagot szeretne elhelyezni.}{Sikertelen művelet, felrobbant aszteroidák raktára nem elérhető. }
\test{Can't place resource in full inventory}{Az aszteroida raktára tele van, a telepes nyersanyagot szeretne elhelyezni benne. }{Sikertelen művelet, a raktárak kapacitása véges, ennek elérése után nem helyezhető bele semmi.}
\test{Can't place resource in unavailable inventory}{Az aszteroida magja nincs teljesen átfúrva, a telepes nyersanyagot szeretne elhelyezni benne.}{Sikertelen művelet, egy nem átfúrt aszteroida raktára nem elérhető.}
\test{SolarFlare}{Napkitörés indul a kör végén.}{Napkitörés hatással van az ott tartózkodó járművekre.}


\section{Tesztelést támogató segéd- és fordítóprogramok specifikálása}
%\comment{Specifikálni kell a tesztelést támogató segédprogramokat. Rövid bemutatással (elvárt funkcionalitás) specifikálni kell a tesztelést támogató segédprogramokat.}
Terveink szerint nem használunk ilyen programokat.


\section{Napló}

\begin{naplo}
	\naplotag{2021.03.27. 16:30 }{ 4 óra }{ Csapat }
	{ 
		Meeting 
	}
    \naplotag{2021.03.28. 18:30 }{ 3 óra }{ Csapat }
	{ 
		Közös feladatmegoldás
	}
	\naplotag{2021.03.28. 19:00 }{ 3 óra }{ Nagy }
	{ 
		Szekvenciák (Craft) + tesztesetek
	}
    \naplotag{2021.03.29. 10:30 }{ 4 óra }{ Dömötör }
	{ 
		Class diagramok
	}
	\naplotag{2021.03.29. 10:30 }{ 1,5 óra }{ Tatai }
	{ 
		Kimenet leírása, dokumentum javítás/kiegészítés
	}
	\naplotag{2021.03.29. 10:30 }{ 1 óra }{ Sike }
	{ 
		UseCase leírások
	}
	\naplotag{2021.03.29. 10:30 }{ 4 óra }{ Dömötör }
	{ 
		UseCase leírások dokumentum összefűzése
	}
\end{naplo}




\end{document}
\clearpage
\section{Napló}

\comment{ 
    A napló tartalmazza az előző beadás óta eltelt időszak történéseit időrendben. A naplóból egyértelműen ki kell derülnie, hogy az egyes anyagrészeket ki és mennyi idő alatt készítette.
    
    A napló bejegyzésekből áll. Minden bejegyzésnek tartalmaznia kell: 
    \begin{itemize}
        \item a történés kezdetének időpontját, nap-óra pontossággal
        \item történés időtartamát, óra felbontással
        \item a szereplő(k) nevét (Kérjük a szereplők VEZETÉKNEVÉT használni)
        \item a tevékenység leírását.
    \end{itemize}
    
    Amennyiben a tevékenységben több szereplő vesz részt, akkor aza tevékenység csak értekezlet lehet, amelynek az eredményei DÖNTÉSEK. A döntéseket precízen meg kell szövegezni (Pl.: Az X objektum Y és Z metódusainak kódját W készíti el Q határidőre).Ha a bejegyzés egyetlen személyhez kötődik, akkor meg kell adni, hogy a tevékenység milyen dologra irányul. A dolog a feladat kapcsán elkészítendő termék, amelynek a (esetleg korábban) beadott anyagban megtalálhatónak kell lenni.A naplóbejegyzés felbontásának egysége szöveges, rajzos anyag esetében az ábra, diagram, vagy kb. fél-egy oldalnyi szöveg. Kódban az egység a metódus. (Pl.: A 3. ábrán látható szekvencia-diagram kidolgozása, vagy az X objektum Y és Z metódusainak kódolása és belövése.
}

\begin{naplo}
    \naplotag{feb. 18. 16h }{ 1 óra }{ Csapat }{ Értekezlet  \newline Döntés: Segédeszközök kiválasztása (git, trello, drive)}
\end{naplo}

\chapter{Részletes tervek}
\comment{A dokumentum célja, hogy pontosan specifikálja az implementálandó osztályokat, beleértve a privát attribútumokat és metódusokat, ezek definícióját is.A dokumentum második fele részletesen be kell mutassa a korábban definiált be-és kimeneti nyelv szintakszisát felhasználva, hogy mely tesztekkel lesz a prototípus ellenőrizve.}
\section{Osztályok és metódusok tervei}
\subsection{Osztály1}
\begin{class-template-responsibility}
    Felelősség leírása
\end{class-template-responsibility}
\begin{class-template-interface}
    Megvalósított interfészek felsorolása
\end{class-template-interface}
\begin{class-template-baseclass}
    Ős-Ősosztály \baseclass Ősosztály... 
\end{class-template-baseclass}
\begin{class-template-attribute}
    \classitem{+A [0..*]}{adattag A}
\end{class-template-attribute}
\comment{Milyen publikus, protected és privát  metódusokkal rendelkezik. Metódusonként precíz leírás, ha szükséges, activity diagram is a metódusban megvalósítandó algoritmusról. Minden olyan metódusnak szerepelnie kell, amelyiket az osztály megvalósít vagy felüldefiniál.}
\begin{class-template-method}
    \classitem{+B(A a) : void}{metódus B}
\end{class-template-method}

\subsection{Osztály2}
\begin{class-template-responsibility}
    Felelősség leírása
\end{class-template-responsibility}
\begin{class-template-statechart}
    A belső működéshez tartozó állapotgép
    \diagram{docs/img/BMElogo}{Az osztály állapotváltozásai}{3cm}
\end{class-template-statechart}

\section{A tesztek részletes tervei, leírásuk a teszt nyelvén}
\comment{A tesztek részletes tervei alatt meg kell adni azokat a bemeneti adatsorozatokat, amelyekkel a program működése ellenőrizhető. Minden bemenő adatsorozathoz definiálni kell, hogy az adatsorozat végrehajtásától a program mely részeinek, funkcióinak ellenőrzését várjuk és konkrétan milyen eredményekre számítunk, ezek az eredmények hogyan vethetők össze a bemenetekkel.A tesztek leírásakor az előző dokumentumban (proto koncepciója) megadott szintakszist kell használni.}

\subsection{Teszteset1}
\begin{test-case-description}
    Szöveges leírás, kb 1-5 mondat
\end{test-case-description}
\begin{test-case-function}
    ...
\end{test-case-function}
\begin{test-case-input}
    A proto nyelvén megadva, pl.:
    \begin{verbatim}
    init world
    0: S 3
    ...
    print
    \end{verbatim}
\end{test-case-input}
\begin{test-case-output}
    A poro kimeneti nyelvén megadva, pl.:
    \begin{verbatim}
    worlddata
    ...
    creatures 0
    \end{verbatim}
\end{test-case-output}

\section{A tesztelést támogató programok tervei}
\comment{A tesztadatok előállítására, a tesztek eredményeinek kiértékelésére szolgáló segédprogramok részletes terveit kell elkészíteni.}
\clearpage
\section{Napló}

\comment{ 
    A napló tartalmazza az előző beadás óta eltelt időszak történéseit időrendben. A naplóból egyértelműen ki kell derülnie, hogy az egyes anyagrészeket ki és mennyi idő alatt készítette.
    
    A napló bejegyzésekből áll. Minden bejegyzésnek tartalmaznia kell: 
    \begin{itemize}
        \item a történés kezdetének időpontját, nap-óra pontossággal
        \item történés időtartamát, óra felbontással
        \item a szereplő(k) nevét (Kérjük a szereplők VEZETÉKNEVÉT használni)
        \item a tevékenység leírását.
    \end{itemize}
    
    Amennyiben a tevékenységben több szereplő vesz részt, akkor aza tevékenység csak értekezlet lehet, amelynek az eredményei DÖNTÉSEK. A döntéseket precízen meg kell szövegezni (Pl.: Az X objektum Y és Z metódusainak kódját W készíti el Q határidőre).Ha a bejegyzés egyetlen személyhez kötődik, akkor meg kell adni, hogy a tevékenység milyen dologra irányul. A dolog a feladat kapcsán elkészítendő termék, amelynek a (esetleg korábban) beadott anyagban megtalálhatónak kell lenni.A naplóbejegyzés felbontásának egysége szöveges, rajzos anyag esetében az ábra, diagram, vagy kb. fél-egy oldalnyi szöveg. Kódban az egység a metódus. (Pl.: A 3. ábrán látható szekvencia-diagram kidolgozása, vagy az X objektum Y és Z metódusainak kódolása és belövése.
}

\begin{naplo}
    \naplotag{feb. 18. 16h }{ 1 óra }{ Csapat }{ Értekezlet  \newline Döntés: Segédeszközök kiválasztása (git, trello, drive)}
\end{naplo}

\setcounter{chapter}{9}
\documentclass[../../projlab]{subfiles}
\begin{document}

\makeatletter

\ifSubfilesClassLoaded{
	\coverpage{10. Leadás}
	%\renewcommand{\filePath}[1]{./../../ #1}
	\def\filePath[#1]{./../../#1}
}{}

\makeatother


\chapter{Prototípus beadása}

\section{Fordítási és futtatási útmutató}
\comment{A feltöltött program fordításával és futtatásával kapcsolatos útmutatás. Ennek tartalmaznia kell leltárszerűen az egyes fájlok pontos nevét, méretét byte-ban, keletkezési idejét, valamint azt, hogy a fájlban mi került megvalósításra.}

\subsection{Fájllista}

\begin{fajllista}
    \fajl
    {Building.java} % Kezdet
    {1 KB} % Méret
    {2021.04.21~12:42~} % Időpont
    {Épületek programosztálya} % Leírás
	\fajl
    {SpaceStation.java} % Kezdet
    {2 KB} % Méret
    {2021.04.21~12:42~} % Időpont
    {Űrállomás programosztálya} % Leírás
	\fajl
    {TeleportGate.java} % Kezdet
    {6 KB} % Méret
    {2021.04.21~12:42~} % Időpont
    {Teleportkapuk programosztálya} % Leírás
	\fajl
    {Robot.java} % Kezdet
    {2 KB} % Méret
    {2021.04.21~12:42~} % Időpont
    {Robotok programosztálya} % Leírás
	\fajl
    {SpaceShip.java} % Kezdet
    {4 KB} % Méret
    {2021.04.21~12:42~} % Időpont
    {Telepesek programosztálya} % Leírás
	\fajl
    {UFO.java} % Kezdet
    {2 KB} % Méret
    {2021.04.21~12:42~} % Időpont
    {UFOk programosztálya} % Leírás
	\fajl
    {Vessel.java} % Kezdet
    {5 KB} % Méret
    {2021.04.21~12:42~} % Időpont
    {Űrjárművek általános programosztálya} % Leírás
	\fajl
    {Asteroid.java} % Kezdet
    {14 KB} % Méret
    {2021.04.21~12:42~} % Időpont
    {Aszteroidák programosztálya} % Leírás
	\fajl
    {Entity.java} % Kezdet
    {2 KB} % Méret
    {2021.04.21~12:42~} % Időpont
    {Entitások általános programosztálya} % Leírás
	\fajl
    {MovableEntity.java} % Kezdet
    {2 KB} % Méret
    {2021.04.21~12:42~} % Időpont
    {Mozgásra képes entitások általános programosztálya} % Leírás
	\fajl
    {AsteroidInventory.java} % Kezdet
    {3 KB} % Méret
    {2021.04.21~12:42~} % Időpont
    {Aszteroidák raktárának programosztálya} % Leírás
	\fajl
    {IInventory.java} % Kezdet
    {2 KB} % Méret
    {2021.04.21~12:42~} % Időpont
    {Raktár interfész} % Leírás
	\fajl
    {InfiniteInventory.java} % Kezdet
    {2 KB} % Méret
    {2021.04.21~12:42~} % Időpont
    {UFOk raktárának programosztálya} % Leírás
	\fajl
    {SSInventory.java} % Kezdet
    {4 KB} % Méret
    {2021.04.21~12:42~} % Időpont
    {Telepesek raktárának programosztálya} % Leírás
	\fajl
    {Coal.java} % Kezdet
    {1 KB} % Méret
    {2021.04.21~12:42~} % Időpont
    {Szén programosztálya} % Leírás
	\fajl
    {Ice.java} % Kezdet
    {2 KB} % Méret
    {2021.04.21~12:42~} % Időpont
    {Vízjég programosztálya} % Leírás
	\fajl
    {Iron.java} % Kezdet
    {1 KB} % Méret
    {2021.04.21~12:42~} % Időpont
    {Vas programosztálya} % Leírás
	\fajl
    {Resource.java} % Kezdet
    {1 KB} % Méret
    {2021.04.21~12:42~} % Időpont
    {Nyersanyagok általános programosztálya} % Leírás
	\fajl
    {Titan.java} % Kezdet
    {1 KB} % Méret
    {2021.04.21~12:42~} % Időpont
    {Titánium programosztálya} % Leírás
	\fajl
    {Uranium.java} % Kezdet
    {2 KB} % Méret
    {2021.04.21~12:42~} % Időpont
    {Uránium programosztálya} % Leírás
	\fajl
    {Item.java} % Kezdet
    {2 KB} % Méret
    {2021.04.21~12:42~} % Időpont
    {Raktárelemek programosztálya} % Leírás
	\fajl
    {TeleportGateItem.java} % Kezdet
    {2 KB} % Méret
    {2021.04.21~12:42~} % Időpont
    {Teleportkapu-elemek programosztálya} % Leírás
	\fajl
    {Recipe.java} % Kezdet
    {5 KB} % Méret
    {2021.04.21~12:42~} % Időpont
    {Receptek általános programosztálya} % Leírás
	\fajl
    {RobotRecipe.java} % Kezdet
    {2 KB} % Méret
    {2021.04.21~12:42~} % Időpont
    {Robot receptjének programosztálya} % Leírás
	\fajl
    {SpaceStationRecipe.java} % Kezdet
    {2 KB} % Méret
    {2021.04.21~12:42~} % Időpont
    {Űrállomás receptjének programosztálya} % Leírás
	\fajl
    {TeleportGateRecipe.java} % Kezdet
    {3 KB} % Méret
    {2021.04.21~12:42~} % Időpont
    {Teleportkapuk receptjének programosztálya} % Leírás
	\fajl
    {Engine.java} % Kezdet
    {3 KB} % Méret
    {2021.04.21~12:42~} % Időpont
    {Játékmotor programosztálya} % Leírás
	\fajl
    {GameManager.java} % Kezdet
    {9 KB} % Méret
    {2021.04.21~12:42~} % Időpont
    {Játékmenedzser programosztálya} % Leírás
	\fajl
    {Main.java} % Kezdet
    {1 KB} % Méret
    {2021.04.21~12:42~} % Időpont
    {Főosztály} % Leírás
	\fajl
    {Player.java} % Kezdet
    {2 KB} % Méret
    {2021.04.21~12:42~} % Időpont
    {Játékosok programosztálya} % Leírás
	\fajl
    {Scene.java} % Kezdet
    {3 KB} % Méret
    {2021.04.21~12:42~} % Időpont
    {Játéktér programosztálya} % Leírás
	\fajl
    {Command.java} % Kezdet
    {1 KB} % Méret
    {2021.04.21~12:42~} % Időpont
    {Általános commandok programosztálya} % Leírás
	\fajl
    {CommandExecutor.java} % Kezdet
    {2 KB} % Méret
    {2021.04.21~12:42~} % Időpont
    {Commandok rendszerezésének, végrehajtásának programosztálya} % Leírás
	\fajl
    {CraftCmd.java} % Kezdet
    {3 KB} % Méret
    {2021.04.21~12:42~} % Időpont
    {Készítés commandjának programosztálya} % Leírás
	\fajl
    {DrillCmd.java} % Kezdet
    {2 KB} % Méret
    {2021.04.21~12:42~} % Időpont
    {Fúrás commandjának programosztálya} % Leírás
	\fajl
    {EndturnCmd.java} % Kezdet
    {1 KB} % Méret
    {2021.04.21~12:42~} % Időpont
    {Kör vége commandjának programosztálya} % Leírás
	\fajl
    {HelpCmd.java} % Kezdet
    {1 KB} % Méret
    {2021.04.21~12:42~} % Időpont
    {Commandok listázásának programosztálya} % Leírás
	\fajl
    {HideCmd.java} % Kezdet
    {2 KB} % Méret
    {2021.04.21~12:42~} % Időpont
    {Elbújás commandjának programosztálya} % Leírás
	\fajl
    {MineCmd.java} % Kezdet
    {2 KB} % Méret
    {2021.04.21~12:42~} % Időpont
    {Bányászat commandjának programosztálya} % Leírás
	\fajl
    {MoveCmd.java} % Kezdet
    {2 KB} % Méret
    {2021.04.21~12:42~} % Időpont
    {Mozgás commandjának programosztálya} % Leírás
	\fajl
    {NeighboursCmd.java} % Kezdet
    {2 KB} % Méret
    {2021.04.21~12:42~} % Időpont
    {Szomszédságok commandjának programosztálya} % Leírás
	\fajl
    {PickupCmd.java} % Kezdet
    {2 KB} % Méret
    {2021.04.21~12:42~} % Időpont
    {Item felszedése commandjának programosztálya} % Leírás
	\fajl
    {PlaceCmd.java} % Kezdet
    {2 KB} % Méret
    {2021.04.21~12:42~} % Időpont
    {Item visszahelyezése commandjának programosztálya} % Leírás
	\fajl
    {SolarDistCmd.java} % Kezdet
    {1 KB} % Méret
    {2021.04.21~12:42~} % Időpont
    {Naptávolság commandjának programosztálya} % Leírás
	\fajl
    {SolarFlareCmd.java} % Kezdet
    {1 KB} % Méret
    {2021.04.21~12:42~} % Időpont
    {Napvihar commandjának programosztálya} % Leírás
	\fajl
    {StartgameCmd.java} % Kezdet
    {2 KB} % Méret
    {2021.04.21~12:42~} % Időpont
    {Játékindítás commandjának programosztálya} % Leírás
	\fajl
    {StatusCmd.java} % Kezdet
    {2 KB} % Méret
    {2021.04.21~12:42~} % Időpont
    {Státuszlekérdezés commandjának programosztálya} % Leírás
	\fajl
    {Namer.java} % Kezdet
    {1 KB} % Méret
    {2021.04.21~12:42~} % Időpont
    {Névgenerálás programosztálya} % Leírás
	\fajl
	{RandomCollection.java} % Kezdet
    {1 KB} % Méret
    {2021.04.21~12:42~} % Időpont
    {Súlyozott lista programosztálya} % Leírás
	\fajl
	{Vector2.java} % Kezdet
    {1 KB} % Méret
    {2021.04.21~12:42~} % Időpont
    {Pozíció típusának programosztálya} % Leírás
	\fajl
	{StationCreator.java} % Kezdet
    {1 KB} % Méret
    {2021.04.21~12:42~} % Időpont
    {Űrállomás létrehozás programosztálya} % Leírás
	\fajl
	{TeleGateCreator.java} % Kezdet
    {2 KB} % Méret
    {2021.04.21~12:42~} % Időpont
    {Teleportkapu létrehozás programosztálya} % Leírás
	\fajl
	{RobotCreator.java} % Kezdet
    {3 KB} % Méret
    {2021.04.21~12:42~} % Időpont
    {Robot létrehozás programosztálya} % Leírás
	\fajl
	{SpaceShipCreator.java} % Kezdet
    {4 KB} % Méret
    {2021.04.21~12:42~} % Időpont
    {Telepes létrehozás programosztálya} % Leírás
	\fajl
	{UfoCreator.java} % Kezdet
    {3 KB} % Méret
    {2021.04.21~12:42~} % Időpont
    {Ufo létrehozás programosztálya} % Leírás
	\fajl
	{AsteroidCreator.java} % Kezdet
    {6 KB} % Méret
    {2021.04.21~12:42~} % Időpont
    {Aszteroida létrehozás programosztálya} % Leírás
	\fajl
	{BuildingCreator.java} % Kezdet
    {3 KB} % Méret
    {2021.04.21~12:42~} % Időpont
    {Épület létrehozás programosztálya} % Leírás
	\fajl
	{MovableEntityCreator.java} % Kezdet
    {2 KB} % Méret
    {2021.04.21~12:42~} % Időpont
    {Mozgásra képes entitás létrehozás programosztálya} % Leírás
	\fajl
	{VesselCreator.java} % Kezdet
    {4 KB} % Méret
    {2021.04.21~12:42~} % Időpont
    {Űrjármű létrehozás programosztálya} % Leírás
	\fajl
	{CoalCreator.java} % Kezdet
    {1 KB} % Méret
    {2021.04.21~12:42~} % Időpont
    {Szén létrehozás programosztálya} % Leírás
	\fajl
	{IceCreator.java} % Kezdet
    {1 KB} % Méret
    {2021.04.21~12:42~} % Időpont
    {Vízjég létrehozás programosztálya} % Leírás
	\fajl
	{IronCreator.java} % Kezdet
    {1 KB} % Méret
    {2021.04.21~12:42~} % Időpont
    {Vas létrehozás programosztálya} % Leírás
	\fajl
	{TeleItemCreator.java} % Kezdet
    {1 KB} % Méret
    {2021.04.21~12:42~} % Időpont
    {Teleportkapu item létrehozás programosztálya} % Leírás
	\fajl
	{TitaniumCreator.java} % Kezdet
    {1 KB} % Méret
    {2021.04.21~12:42~} % Időpont
    {Titán létrehozás programosztálya} % Leírás
	\fajl
	{UraniumCreator.java} % Kezdet
    {1 KB} % Méret
    {2021.04.21~12:42~} % Időpont
    {Urán létrehozás programosztálya} % Leírás
	\fajl
	{EntityCreator.java} % Kezdet
    {1 KB} % Méret
    {2021.04.21~12:42~} % Időpont
    {Entitás létrehozás programosztálya} % Leírás
	\fajl
	{FileOpener.java} % Kezdet
    {8 KB} % Méret
    {2021.04.21~12:42~} % Időpont
    {Fájl megnyitás programosztálya} % Leírás
	\fajl
	{GMCreator.java} % Kezdet
    {6 KB} % Méret
    {2021.04.21~12:42~} % Időpont
    {GameManager létrehozás programosztálya} % Leírás
	\fajl
	{ItemCreator.java} % Kezdet
    {3 KB} % Méret
    {2021.04.21~12:42~} % Időpont
    {Item létrehozás programosztálya} % Leírás
	\fajl
	{MrConnector.java} % Kezdet
    {2 KB} % Méret
    {2021.04.21~12:42~} % Időpont
    {Aszteroidakapcsolatok építésének programosztálya} % Leírás
	\fajl
	{PlayerCreator.java} % Kezdet
    {2 KB} % Méret
    {2021.04.21~12:42~} % Időpont
    {Játékos létrehozás programosztálya} % Leírás
	\fajl
	{SceneCreator.java} % Kezdet
    {1 KB} % Méret
    {2021.04.21~12:42~} % Időpont
    {Játéktér létrehozás programosztálya} % Leírás
\end{fajllista}

\subsection{Fordítás}
%\comment{A fenti listában szereplő forrásfájlokból milyen műveletekkel lehet a bináris, futtatható kódot előállítani. Az előállításhoz csak a 2. Követelmények c. dokumentumban leírt környezetet szabad előírni.}
A fenti fájlokból a javac parancs kiadásával fordítható a tesztelőprogram (az alábbi módon).
Esetleg GUI-val rendelkező fejlesztőkörnyezetből, az annak megfelelő módon.

\begin{verbatim}
    javac -d bin *.java
\end{verbatim}

\subsection{Futtatás}
%\comment{A futtatható kód elindításával kapcsolatos teendők leírása. Az indításhoz csak a 2. Követelmények c. dokumentumban leírt környezetet szabad előírni.}

Futtatáshoz a következő parancsokat kell kiadni.
\begin{verbatim}
    cd bin
    java SkeletonEntry
\end{verbatim}

\section{Tesztek jegyzőkönyvei}

\subsection{Teszteset1}
\comment{Az alábbi táblázatot a megismételt (hibás) tesztek esetén kell kitölteni minden ismétléshez egyszer. Ha szükséges, akkor a valós kimenet is mellékelhető mint a teszt eredménye}
\testFAIL{????}{Ápr. 25. 13:50}{ A teszt eredménye}{A teszt hibás futásának okai}{Változtatások}
\comment{Az alábbi táblázatot az utolsó, sikeres tesztfuttatáshoz kell kitölteni}
\testOK{????}{Ápr. 25. 14:00}



\section{Értékelés}
\comment{A projekt kezdete óta az értékelésig eltelt időben tagokra bontva, százalékban.}
\begin{ertekeles}
    \ertekelestag{Sike Ádám}{ E8Z277 }{18\%}
    \ertekelestag{Dömötör Péter}{ G2Y5TI }{28\%}
    \ertekelestag{Gao Tong}{ I2SVOS }{18\%}
    \ertekelestag{Nagy Beáta}{ GPOGC5 }{18\%}
    \ertekelestag{Tatai Titusz Miklós}{ IJHLYX }{18\%}
\end{ertekeles}


\section{Napló}

\begin{naplo}

    \naplotag{2021.04.08. 16:00 }{ 4 óra }{ Gao Tong }
	{ 
		Az előző leadásnál lemaradt - Osztályleírások
	}
    \naplotag{2021.04.09. 16:00 }{ 4 óra }{ Gao Tong}
	{ 
		Az előző leadásnál lemaradt - Osztályleírások
	}
    \naplotag{2021.04.14. 10:15 }{ 10 perc }{ Csapat }
	{ 
		Rövid meeting 
	}
     \naplotag{2021.04.19. 21:00 }{ 2 óra }{ Dömötör Péter }
	{ 
		Refaktorozás
	}
    \naplotag{2021.04.16. 18:00 }{ 2 óra }{ Tatai Titusz }
	{ 
		Fájl beolvasás - alapok
	}
	\naplotag{2021.04.17. 14:15 }{ 10 perc }{ Csapat }
	{ 
		Rövid meeting 
	}
    \naplotag{2021.04.17. 14:00 }{ 6 óra }{ Nagy Beáta }
	{ 
		Receptek
	}
    \naplotag{2021.04.17. 12:00 }{ 4 óra }{ Sike Ádám }
	{ 
		Fájl beolvasás - szövegfeldolgozás
	}
     \naplotag{2021.04.17. 15:00 }{ 6 óra }{Dömötör Péter }
	{ 
		Parancsok létrehozása, compile, világgenerálás
	}
    \naplotag{2021.04.17. 16:00 }{ 4 óra }{ Tatai Titusz }
	{ 
		Ufók működése
	}
    \naplotag{2021.04.17. 17:00 }{ 3 óra }{ Gao Tong}
	{ 
		Apróbb munkálatok
	}
    \naplotag{2021.04.18. 18:00 }{ 2 óra }{ Sike Ádám }
	{ 
		Fájl beolvasás - fájl levelei
	}
     \naplotag{2021.04.18. 20:00 }{ 2 óra }{ Nagy Beáta }
	{ 
		Entitás implementáció
	}
    \naplotag{2021.04.19. 20:00 }{ 2 óra }{ Sike Ádám }
	{ 
		Fájl beolvasás - Scene létrehozása
	}
    \naplotag{2021.04.20. 10:00 }{ 4 óra }{ Csapat }
	{ 
		Befejező lépések
	}

\end{naplo}

\end{document}
\clearpage
\section{Napló}

\comment{ 
    A napló tartalmazza az előző beadás óta eltelt időszak történéseit időrendben. A naplóból egyértelműen ki kell derülnie, hogy az egyes anyagrészeket ki és mennyi idő alatt készítette.
    
    A napló bejegyzésekből áll. Minden bejegyzésnek tartalmaznia kell: 
    \begin{itemize}
        \item a történés kezdetének időpontját, nap-óra pontossággal
        \item történés időtartamát, óra felbontással
        \item a szereplő(k) nevét (Kérjük a szereplők VEZETÉKNEVÉT használni)
        \item a tevékenység leírását.
    \end{itemize}
    
    Amennyiben a tevékenységben több szereplő vesz részt, akkor aza tevékenység csak értekezlet lehet, amelynek az eredményei DÖNTÉSEK. A döntéseket precízen meg kell szövegezni (Pl.: Az X objektum Y és Z metódusainak kódját W készíti el Q határidőre).Ha a bejegyzés egyetlen személyhez kötődik, akkor meg kell adni, hogy a tevékenység milyen dologra irányul. A dolog a feladat kapcsán elkészítendő termék, amelynek a (esetleg korábban) beadott anyagban megtalálhatónak kell lenni.A naplóbejegyzés felbontásának egysége szöveges, rajzos anyag esetében az ábra, diagram, vagy kb. fél-egy oldalnyi szöveg. Kódban az egység a metódus. (Pl.: A 3. ábrán látható szekvencia-diagram kidolgozása, vagy az X objektum Y és Z metódusainak kódolása és belövése.
}

\begin{naplo}
    \naplotag{feb. 18. 16h }{ 1 óra }{ Csapat }{ Értekezlet  \newline Döntés: Segédeszközök kiválasztása (git, trello, drive)}
\end{naplo}

\documentclass[../../projlab]{subfiles}
\begin{document}

\makeatletter

\ifSubfilesClassLoaded{
	\coverpage{10. Leadás}
	%\renewcommand{\filePath}[1]{./../../ #1}
	\def\filePath[#1]{./../../#1}
}{}

\makeatother

\chapter{Grafikus felület specifikációja}

\section{A grafikus interfész}
%\comment{A menürendszer, a kezelői felület grafikus képe. A grafikus felület megjelenését, a használt ikonokat, stb screenshot-szerű képekkel kell bemutatni. Az építészetben ez a homlokzati terv.}
\diagram{docs/bead11/img/InGame.png}{Példa kép}{15cm}

\section{A grafikus rendszer architektúrája}
%\comment{A felület működésének elve, a grafikus rendszer architektúrája (struktúra diagramok). A struktúra diagramokon a prototípus azon és csak azon osztályainak is szerepelnie kell, amelyekhez a grafikus felületet létrehozó osztályok kapcsolódnak.}

\subsection{A felület működési elve}
\comment{Le kell írni, hogy a grafikai megjelenésért felelős osztályok, objektumok hogyan kapcsolódnak a meglevő rendszerhez, a megjelenítés során mi volt az alapelv. Törekedni kell az MVC megvalósításra. Alapelvek lehetnek: \textbf{push} alapú: a modell értesíti a felületet, hogy változott; \textbf{pull} alapú: a felület kérdezi le a modellt, hogy változott-e; \textbf{kevert}: a kettő kombinációja.}

\subsection{A felület osztály-struktúrája}
\begin{figure}[H] 
    \centering 
    \includegraphics[width=1\textwidth]{docs/bead11/class_diag} 
    \caption{Class diagram} 
\end{figure} 

\section{A grafikus objektumok felsorolása}
%\comment{Az új osztályok felsorolása. Az régi osztályok közül azoknak a felsorolása, ahol változás volt. Ezek esetén csak a változásokat kell leírni.}

\subsection{Új osztályok}

\subsubsection{ActionBar}
\begin{class-template-responsibility}
    Ha aktív játékmenet van, ez a panel gyűjti össze az aktuálisan kiválasztott entitáson végrehajtható műveleteket.  
\end{class-template-responsibility}
\begin{class-template-interface}
    $\emptyset$
\end{class-template-interface}
\begin{class-template-baseclass}
    JPanel
\end{class-template-baseclass}
\begin{class-template-attribute}
    \item[] $\emptyset$
\end{class-template-attribute} 
\begin{class-template-method}
    \classitem{+ActionBar()}{Konstruktor. A panel megjelenítési tulajdonságait beállítja. Lekérdezi a kiválasztott entitást (getSelectedEntity), és a hozzá tartozó lehetséges parancsokat (getActions), minden parancsnak egy gombot létrehoz a felületen. A gombokhoz action listenereket rendel hozzá, akik a parancsok tényleges végrehajtásáért felelnek, meghívják a megfelelő metódust (doAction). }
\end{class-template-method}

\subsubsection{ConsoleLogPanel}
\begin{class-template-responsibility}
    A játékmenet konzolos logolásáért felelős panel.  
\end{class-template-responsibility}
\begin{class-template-interface}
    StateChangedListener
\end{class-template-interface}
\begin{class-template-baseclass}
    JPanel
\end{class-template-baseclass}
\begin{class-template-attribute}
    \classitem{-log:String[0..*]}{Az aktuális log (konzolos műveletek stringekben tárolva).}
\end{class-template-attribute}
\begin{class-template-method}
    \classitem{+ConsoleLogPanel()}{Konstruktor. A panel megjelenítési tulajdonságait beállítja.}
    \classitem{+stateChanged()}{}
\end{class-template-method}

\subsubsection{InGameGUI}
\begin{class-template-responsibility}
    Ha aktív játékmenet van, ez a panel gyűjti össze a játékot kirajzoló ablak komponenseit.  
\end{class-template-responsibility}
\begin{class-template-interface}
    $\emptyset$
\end{class-template-interface}
\begin{class-template-baseclass}
    JPanel
\end{class-template-baseclass}
\begin{class-template-attribute}
    \item[] $\emptyset$
\end{class-template-attribute} 
\begin{class-template-method}
    \classitem{+InGameGUI()}{Konstruktor. A panel megjelenítési tulajdonságait beállítja. Létrehoz egy SelectionInspectorPanel, MapViewPanel, ActionBar, ConsoleLogPanel, PlayerList objektumot, és elhelyezi ezeket megfelelően az ablakon.}
\end{class-template-method}

\subsubsection{INotifyPropertyChanged}
\begin{class-template-responsibility}
    Interfész: a modellnek azok az elemei valósítják meg, amik értesítik a viewt a modellben létrejövő változásokról. A rájuk feliratkozó listener osztályok majd ennek megfelelően változtatják a view különböző elemeit.
\end{class-template-responsibility}
\begin{class-template-interface}
    $\emptyset$
\end{class-template-interface}
\begin{class-template-baseclass}
    $\emptyset$
\end{class-template-baseclass}
\begin{class-template-attribute}
    \item[] $\emptyset$
\end{class-template-attribute}
\begin{class-template-method}
    \classitem{+addListener(l:StateChangedListener)}{Rá feliratkozó listener osztály hozzáadásáért felel. }
    \classitem{+removeListener(l:StateChangedListener)}{Rá feliratkozó listener osztály eltávolításáért felel. }
\end{class-template-method}

\subsubsection{MainMenu}
\begin{class-template-responsibility}
    A játék indulása előtt ez a panel jelenik meg, itt lehet ténylegesen elindítani egy játékmenetet.
\end{class-template-responsibility}
\begin{class-template-interface}
    $\emptyset$
\end{class-template-interface}
\begin{class-template-baseclass}
    JPanel
\end{class-template-baseclass}
\begin{class-template-attribute}
    \item[] $\emptyset$
\end{class-template-attribute}
\begin{class-template-method}
    \classitem{+MainMenu()}{Konstruktor. A panel megjelenítési tulajdonságait beállítja. Létrehoz egy "new game" gombot, amihez egy action listenert rendel, aki egy új játékmenet elindításáért felelős: az enginen keresztül beállítja a játékosok számát (setPlayerCount), beceneveiket, és elindít egy új játékot (StartGame). }
\end{class-template-method}

\subsubsection{MainWindow}
\begin{class-template-responsibility}
    A programot kirajzoló főablak. 
\end{class-template-responsibility}
\begin{class-template-interface}
    StateChangedListener
\end{class-template-interface}
\begin{class-template-baseclass}
    JFrame
\end{class-template-baseclass}
\begin{class-template-attribute}
    \classitem{-eng:Engine}{Az aktuális játékfolyamot tárolja.}
    \classitem{-active:JPanel}{A jelenleg aktív panel, a játék állapotától függ.}
\end{class-template-attribute}
\begin{class-template-method}
    \classitem{+MainWindow()}{Konstruktor. Az ablak megjelenítési tulajdonságait beállítja. Inicializálja a scene és active változókat (az active alapértelmezetten egy újonnan létrehozott MainMenu, hiszen menüvel kezdődik a játék). Az engine-hez hozzáadja magát, mint listenert, hogy értesüljön a játék állapotának változásairól.}
    \classitem{+stateChanged()}{Modell állapotváltozása esetén hívódik meg. Lekérdezi, hogy az engine állapota mire változott, és ettől függően hajt végre műveletet. MainMenu állapotnál létrehoz és aktívvá tesz egy új MainMenu objektumot, ezt fogja kirajzolni az ablakban. InGame esetén létrehoz és aktívvá tesz egy új InGameGUI objektumot, ezt fogja kirajzolni. Won és Lost esetén vége a játéknak, nem történik változás.}
\end{class-template-method}

\subsubsection{MapViewPanel}
\begin{class-template-responsibility}
    Az aktív játék aktuális állapotát kirajzoló osztály.  
\end{class-template-responsibility}
\begin{class-template-interface}
    StateChangedListener
\end{class-template-interface}
\begin{class-template-baseclass}
    JPanel
\end{class-template-baseclass}
\begin{class-template-attribute}
    \classitem{-scene:Scene}{Az entitásokat összefogó osztályra hivatkozó attribútum.}
\end{class-template-attribute}
\begin{class-template-method}
    \classitem{+MapViewPanel()}{Konstruktor. A panel megjelenítési tulajdonságait beállítja. Inicializálja a scene attribútumot. Az enginen keresztül elért GameManager-hez hozzáadja magát, mint listenert, hogy értesüljön az aktuális entitás megváltozásáról.}
    \classitem{+paint(g:Graphics)}{Az ablakot kirajzolja grafikus felületre. Az ősét kirajzolja, beállítja a hátteret, kirajzolja az összes entitást a megadott színnel és mérettel, és az aszteroidák közti kapcsolatokat vonalakkal. Az aktuálisan kijelölt entitást kiemeli külön színnel.}
    \classitem{+stateChanged()}{Újrarajzolja az ablakot, hogy az aktuálisan kijelölt entitás színe megváltozzon, és a játékmenet modelljének egyéb változásait kövesse a view is.}
\end{class-template-method}

\subsubsection{PlayerList}
\begin{class-template-responsibility}
    A játékosok neveinek felsorolásáért, az aktuális játékos kiemeléséért felelős panel.  
\end{class-template-responsibility}
\begin{class-template-interface}
    StateChangedListener
\end{class-template-interface}
\begin{class-template-baseclass}
    JPanel
\end{class-template-baseclass}
\begin{class-template-attribute}
    \item[] $\emptyset$
\end{class-template-attribute}
\begin{class-template-method}
    \classitem{+PlayerList()}{Konstruktor. A panel megjelenítési tulajdonságait beállítja, meghívja a doPlayerList metódust.}
    \classitem{+doPlayerList()}{Lekérdezi a játék managert, majd tőle az összes játékost. Mindegyikhez létrehoz egy új labelt, és a címkét beállítja a játékos nevére, így megjelenik a játékosok listája.}
    \classitem{+stateChanged()}{}
\end{class-template-method}

\subsubsection{SelectionInspectorPanel}
\begin{class-template-responsibility}
    Az entitások listáját kirajzoló, és aktuális entitást grafikusan kiválasztó panel.  
\end{class-template-responsibility}
\begin{class-template-interface}
    StateChangedListener
\end{class-template-interface}
\begin{class-template-baseclass}
    JPanel
\end{class-template-baseclass}
\begin{class-template-attribute}
    \classitem{-scene:Scene}{Az entitásokat összefogó osztályra hivatkozó attribútum.}
    \classitem{-selected:Entity}{A jelenleg kiválasztott entitás}
    \classitem{-t2:JTextArea}{A kiválasztott entitás részletesebb státuszát tárolja szövegesen.}
\end{class-template-attribute}
\begin{class-template-method}
    \classitem{+SelectionInspectorPanel()}{Konstruktor. A panel megjelenítési tulajdonságait beállítja. Inicializálja a scene attribútumot. Az enginen keresztül elért GameManager-hez hozzáadja magát, mint listenert, hogy értesüljön az aktuális entitás megváltozásáról. A scene-től lekérdezi az összes entitást, amiket egy listában eltárol (JList), ennek az elemeit rajzolja ki. A listához hozzáad egy listenert, akit a kijelölt elem változásáról értesít. A listener meghívja a selectionChanged metódust, átadva a kijelölt elem indexét.}
    \classitem{+selectionChanged(index:int)}{A kijelölt entitás megváltozásakor hívódik meg. A scene-től lekérdezi a paraméterként megkapott indexhez tartozó entitást, majd beállítja ezt aktuális entitásnak (setSelectedEntity).}
    \classitem{+stateChanged()}{Ha a kijelölt entitás megváltozik, lekérdezi a részletes információit (printStatus metódusa), és ezzel frissíti a t2 attribútum tartalmát.}
\end{class-template-method}

\subsubsection{StateChangedListener}
\begin{class-template-responsibility}
    Interfész: A modell állapotváltozásaira feliratkozó listenerek valósítják meg. 
\end{class-template-responsibility}
\begin{class-template-interface}
    $\emptyset$
\end{class-template-interface}
\begin{class-template-baseclass}
    $\emptyset$
\end{class-template-baseclass}
\begin{class-template-attribute}
    \item[] $\emptyset$
\end{class-template-attribute}
\begin{class-template-method}
    \classitem{+stateChanged()}{Az állapotváltozás hatására az osztályban végbement változásokat hajtja végre.}
\end{class-template-method}

\subsection{Megváltozott osztályok}

\subsubsection{Engine}
\begin{class-template-responsibility}
    A program indításáért, egy játékfolyam bezárásáért felelős osztály.
\end{class-template-responsibility}
\begin{class-template-interface}
    INotifyPropertyChanged 
\end{class-template-interface}
\begin{class-template-baseclass}
    $\emptyset$
\end{class-template-baseclass}
\begin{class-template-attribute}
    \classitem{-gameState:State}{A játék aktuális állapotát jelzi egy enumban. Lehetséges értékei: MainMenu (játék még nem indult el, menüben vagyunk), InGame (a játék jelenleg folyik), Won (a játéknak vége, nyertek a játékosok). Lost(a játéknak vége, vesztettek a játékosok)}
    \classitem{-listeners:StateChangedListener[0..*] }{Az osztályok listája, akik fel vannak iratkozva az Engine eseményeire. Őket értesíti, ha megváltozik a játék állapota.}
\end{class-template-attribute}
\begin{class-template-method}
    \classitem{+getGameState():State}{Aktuális játékállapot gettere}
    \classitem{+setGameState(gameState:State)}{Aktuális játékállapot settere. Értesíti a listenereket az állapotváltozásról, ha ez bekövetkezik. }
    \classitem{+addListener(l:StateChangedListener) }{Feliratkozás az Engine eseményeire. A paraméterként megadott osztály feliratkozik: hozzáadja a listeners listához az elemet.}
    \classitem{+removeListener(l:StateChangedListener) }{Leiratkozás az Engine eseményeiről. A paraméterként megadott osztály leiratkozik: eltávolítja a listeners listából az elemet.}
    \classitem{+notifyListeners()}{Értesíti az összes listenert a listán, hogy változás történt a modellben. Végigiterál a listenerek listáján, mindenkinek meghívja a stateChanged metódusát.}
\end{class-template-method}

\subsubsection{Entity}
\begin{class-template-responsibility}
    Egy adott entitás (vizuális megjelenítéssel rendelkező játékelem) osztálya. 
\end{class-template-responsibility}
\begin{class-template-interface}
    $\emptyset$
\end{class-template-interface}
\begin{class-template-baseclass}
    $\emptyset$
\end{class-template-baseclass}
\begin{class-template-attribute}
    \item[] $\emptyset$
\end{class-template-attribute} 
\begin{class-template-method}
    \classitem{+printStatus():String}{Egy adott entitáshoz tartozó státuszinformációk. A leszármazottjai felüldefiniálják a náluk releváns információkkal. Pl. egy aszteroida kiírja a kéreg mélységét, a rajta állomásozó járművek listáját, a benne található itemeket, rajta lévő teleportkapukat stb. Gyakorlatilag csak stringeket illesztünk benne össze információkiíráshoz, ezért nem szükséges a működéséről szekvenciadiagramot készíteni. Az átláthatóság miatt az összes leszármazottat nem jelöltük aki megvalósítja a metódust, mindegyiknél a számára releváns információkkal definiálódik felül a metódus, más stringet ad vissza. }
    \classitem{+getActions():List<String>}{Visszaadja az actions attribútum értékét azoknál a leszármazottaknál, ahol ez releváns, máshol null-t ad vissza, jelezve, hogy nincs ilyen attribútum.}
    \classitem{+doAction(name:String)}{A paraméterként megadott Stringet megkeresi az actions listában, és ha megtalálja, az értékének megfelelő műveletet hajt végre. Az ActionBar osztályban kijelölt műveletek így hajthatók végre.}
\end{class-template-method}

\subsubsection{GameManager}
\begin{class-template-responsibility}
    A játék menetének irányításáért felelős osztály. 
\end{class-template-responsibility}
\begin{class-template-interface}
    INotifyPropertyChanged 
\end{class-template-interface}
\begin{class-template-baseclass}
    $\emptyset$
\end{class-template-baseclass}
\begin{class-template-attribute}
    \classitem{-selectedEntity:Entity}{Aktuálisan kiválasztott entitás}
    \classitem{-listeners:StateChangedListener[0..*] }{Az osztályok listája, akik fel vannak iratkozva a GameManager eseményeire. Ha változik a kiválasztott játékos, ők fognak értesítést kapni erről.}
\end{class-template-attribute}
\begin{class-template-method}
    \classitem{+getSelectedEntity():Entity}{Aktuálisan kiválasztott entitás gettere}
    \classitem{+setSelectedEntity(e:Entity)}{Aktuálisan kiválasztott entitás settere. Értesíti a listenereket a kiválasztott entitás megváltozásáról, ha ez ténylegesen változik. }
    \classitem{+addListener(l:StateChangedListener) }{Feliratkozás a GameManager eseményeire. A paraméterként megadott osztály feliratkozik: hozzáadja a listeners listához az elemet.}
    \classitem{+removeListener(l:StateChangedListener) }{Leiratkozás a GameManager eseményeiről. A paraméterként megadott osztály leiratkozik: eltávolítja a listeners listából az elemet.}
    \classitem{+notifyListeners()}{Értesíti az összes listenert a listán, hogy változás történt a modellben. Végigiterál a listenerek listáján, mindenkinek meghívja a stateChanged metódusát.}
\end{class-template-method}

\subsubsection{Robot}
\begin{class-template-responsibility}
    Egy robot osztálya. 
\end{class-template-responsibility}
\begin{class-template-interface}
    $\emptyset$
\end{class-template-interface}
\begin{class-template-baseclass}
    Entity \baseclass MovableEntity \baseclass Vessel
\end{class-template-baseclass}
\begin{class-template-attribute}
    \classitem{-actions:String[0..*]}{A robottal végrehajtható műveletek listája (move, hide, drill).}
\end{class-template-attribute}
\begin{class-template-method}
    \classitem{+getActions():List<String>}{Visszaadja az actions attribútum értékét.}
    \classitem{+doAction(name:String)}{A paraméterként megadott Stringet megkeresi az actions listában, és ha megtalálja, az értékének megfelelő műveletet hajt végre. Gyakorlatilag meghívja a move, hide vagy drill metódust. }
\end{class-template-method}

\subsubsection{SpaceShip}
\begin{class-template-responsibility}
    Egy telepes osztálya. 
\end{class-template-responsibility}
\begin{class-template-interface}
    $\emptyset$
\end{class-template-interface}
\begin{class-template-baseclass}
    Entity \baseclass MovableEntity \baseclass Vessel
\end{class-template-baseclass}
\begin{class-template-attribute}
    \classitem{-actions:String[0..*]}{A telepessel végrehajtható műveletek listája (move, hide, drill, mine, activate, create, stash).}
\end{class-template-attribute}
\begin{class-template-method}
    \classitem{+getActions():List<String>}{Visszaadja az actions attribútum értékét.}
    \classitem{+doAction(name:String)}{A paraméterként megadott Stringet megkeresi az actions listában, és ha megtalálja, az értékének megfelelő műveletet hajt végre. Gyakorlatilag meghívja a hozzá tartozó metódust, és megadja az esetleges paramétereket. }
\end{class-template-method}

\subsubsection{UFO}
\begin{class-template-responsibility}
    Egy UFO osztálya. 
\end{class-template-responsibility}
\begin{class-template-interface}
    $\emptyset$
\end{class-template-interface}
\begin{class-template-baseclass}
    Entity \baseclass MovableEntity \baseclass Vessel
\end{class-template-baseclass}
\begin{class-template-attribute}
    \classitem{-actions:String[0..*]}{Az UFOval végrehajtható műveletek listája (move, mine).}
\end{class-template-attribute}
\begin{class-template-method}
    \classitem{+getActions():List<String>}{Visszaadja az actions attribútum értékét.}
    \classitem{+doAction(name:String)}{A paraméterként megadott Stringet megkeresi az actions listában, és ha megtalálja, az értékének megfelelő műveletet hajt végre. Gyakorlatilag meghívja a move vagy mine metódust. }
\end{class-template-method}

\section{Kapcsolat az alkalmazói rendszerrel}
%\comment{Szekvencia-diagramokon ábrázolni kell a grafikus rendszer működését. Konzisztens kell legyen az előző alfejezetekkel. Minden metódus, ami ott szerepel, fel kell tűnjön valamelyik szekvenciában. Minden metódusnak, ami szekvenciában szerepel, szereplnie kell a valamelyik osztálydiagramon.}
\begin{figure}[H] 
    \centering 
    \includegraphics[width=1\textwidth]{docs/bead11/seq/addlistener} 
    \caption{Az addListener-t az interfészen hívjuk meg, mert a két megvalósításnál (GameManager,Engine) ugyanez a műveletsor megy végbe, így két külön diagram nem készült róla.} 
\end{figure} 

\begin{figure}[H] 
    \centering 
    \includegraphics[width=1\textwidth]{docs/bead11/seq/doPlayerList} 
    \caption{Class diagram} 
\end{figure} 


\begin{figure}[H] 
    \centering 
    \includegraphics[width=1\textwidth]{docs/bead11/seq/notifylistenersEngine} 
    \caption{Class diagram} 
\end{figure} 

\begin{figure}[H] 
    \centering 
    \includegraphics[width=1\textwidth]{docs/bead11/seq/notifylistenersManager} 
    \caption{Class diagram} 
\end{figure} 

\begin{figure}[H] 
    \centering 
    \includegraphics[width=1\textwidth]{docs/bead11/seq/paint} 
    \caption{Class diagram} 
\end{figure} 

\begin{figure}[H] 
    \centering 
    \includegraphics[width=1\textwidth]{docs/bead11/seq/removelistener} 
    \caption{A removeListener-t az interfészen hívjuk meg, mert a két megvalósításnál (GameManager,Engine) ugyanez a műveletsor megy végbe, így két külön diagram nem készült róla.} 
\end{figure}

\begin{figure}[H] 
    \centering 
    \includegraphics[width=1\textwidth]{docs/bead11/seq/selectionChanged} 
    \caption{Class diagram} 
\end{figure}

\begin{figure}[H] 
    \centering 
    \includegraphics[width=1\textwidth]{docs/bead11/seq/setEntity} 
    \caption{Class diagram} 
\end{figure} 

\begin{figure}[H] 
    \centering 
    \includegraphics[width=1\textwidth]{docs/bead11/seq/setGameState} 
    \caption{Class diagram} 
\end{figure} 

\begin{figure}[H] 
    \centering 
    \includegraphics[width=1\textwidth]{docs/bead11/seq/stateChanged} 
    \caption{Class diagram} 
\end{figure}

\begin{figure}[H] 
    \centering 
    \includegraphics[width=0.5\textwidth]{docs/bead11/seq/stateChanged2} 
    \caption{Class diagram} 
\end{figure}

\begin{figure}[H] 
    \centering 
    \includegraphics[width=1\textwidth]{docs/bead11/seq/stateChanged3} 
    \caption{Class diagram} 
\end{figure}


\newpage
\section{Napló}

\begin{naplo}
    \naplotag{2021.04.24. 20:00}{ 1 óra }{ Csapat }{
        GUI tervezés
         }
    \naplotag{2021.04.25. 16:00}{ 2 óra }{ Dömötör Péter }{
        Grafikus interfész prototípus
         }
    \naplotag{2021.04.25. 17:00}{ 1 óra }{ Sike Ádám }{
        Szekvencia diagramok
         }
    \naplotag{2021.04.25. 18:00}{ 1 óra }{ Gao Tong}{
        Szekvencia diagramok
         }
    \naplotag{2021.04.25. 18:00}{ 1 óra }{ Tatai Titusz }{
        Grafikus interfész  javítása
         }
    \naplotag{2021.04.25. 18:00}{ 1 óra }{ Nagy Beáta }{
        Szekvencia diagramok
         }
    \naplotag{2021.04.25. 19:00}{ 1 óra }{ Csapat }{
        Apróságok és pár dolog finomítása
         }

    \naplotag{2021.04.26. 07:00}{2 óra}{Nagy Beáta}{
        Osztályleírások
        }
       
\end{naplo}

\end{document}
\clearpage
\section{Napló}

\comment{ 
    A napló tartalmazza az előző beadás óta eltelt időszak történéseit időrendben. A naplóból egyértelműen ki kell derülnie, hogy az egyes anyagrészeket ki és mennyi idő alatt készítette.
    
    A napló bejegyzésekből áll. Minden bejegyzésnek tartalmaznia kell: 
    \begin{itemize}
        \item a történés kezdetének időpontját, nap-óra pontossággal
        \item történés időtartamát, óra felbontással
        \item a szereplő(k) nevét (Kérjük a szereplők VEZETÉKNEVÉT használni)
        \item a tevékenység leírását.
    \end{itemize}
    
    Amennyiben a tevékenységben több szereplő vesz részt, akkor aza tevékenység csak értekezlet lehet, amelynek az eredményei DÖNTÉSEK. A döntéseket precízen meg kell szövegezni (Pl.: Az X objektum Y és Z metódusainak kódját W készíti el Q határidőre).Ha a bejegyzés egyetlen személyhez kötődik, akkor meg kell adni, hogy a tevékenység milyen dologra irányul. A dolog a feladat kapcsán elkészítendő termék, amelynek a (esetleg korábban) beadott anyagban megtalálhatónak kell lenni.A naplóbejegyzés felbontásának egysége szöveges, rajzos anyag esetében az ábra, diagram, vagy kb. fél-egy oldalnyi szöveg. Kódban az egység a metódus. (Pl.: A 3. ábrán látható szekvencia-diagram kidolgozása, vagy az X objektum Y és Z metódusainak kódolása és belövése.
}

\begin{naplo}
    \naplotag{feb. 18. 16h }{ 1 óra }{ Csapat }{ Értekezlet  \newline Döntés: Segédeszközök kiválasztása (git, trello, drive)}
\end{naplo}

\setcounter{chapter}{12}
\chapter{Grafikus felület specifikációja}

\section{Fordítási és futtatási útmutató}
\comment{A feltöltött program fordításával és futtatásával kapcsolatos útmutatás. Ennek tartalmaznia kell leltárszerűen az egyes fájlok pontos nevét, méretét byte-ban, keletkezési idejét, valamint azt, hogy a fájlban mi került megvalósításra.}

\subsection{Fájllista}

\begin{fajllista}
    
\end{fajllista}

\subsection{Fordítás}
\comment{A fenti listában szereplő forrásfájlokból milyen műveletekkel lehet a bináris, futtatható kódot előállítani. Az előállításhoz csak a 2. Követelmények c. dokumentumban leírt környezetet szabad előírni.}

\begin{verbatim}
javac -d bin *.java
\end{verbatim}

\subsection{Futtatás}
\comment{A futtatható kód elindításával kapcsolatos teendők leírása. Az indításhoz csak a 2. Követelmények c. dokumentumban leírt környezetet szabad előírni.}

\begin{verbatim}
cd bin
java Main.java
\end{verbatim}

\clearpage
\section{Értékelés}
\comment{A projekt kezdete óta az értékelésig eltelt időben tagokra bontva, százalékban.}
\begin{ertekeles}
    \ertekelestag{Dömötör Péter}{ G2Y5TI }{\%}
    \ertekelestag{Sike Ádám}{ E8Z277 }{\%}
    \ertekelestag{Gao Tong}{ I2SVOS }{\%}
    \ertekelestag{Nagy Beáta}{ GPOGC5 }{\%}
    \ertekelestag{Tatai Titusz Miklós}{ IJHLYX }{\%}
\end{ertekeles}


\clearpage
\section{Napló}

\comment{ 
    A napló tartalmazza az előző beadás óta eltelt időszak történéseit időrendben. A naplóból egyértelműen ki kell derülnie, hogy az egyes anyagrészeket ki és mennyi idő alatt készítette.
    
    A napló bejegyzésekből áll. Minden bejegyzésnek tartalmaznia kell: 
    \begin{itemize}
        \item a történés kezdetének időpontját, nap-óra pontossággal
        \item történés időtartamát, óra felbontással
        \item a szereplő(k) nevét (Kérjük a szereplők VEZETÉKNEVÉT használni)
        \item a tevékenység leírását.
    \end{itemize}
    
    Amennyiben a tevékenységben több szereplő vesz részt, akkor aza tevékenység csak értekezlet lehet, amelynek az eredményei DÖNTÉSEK. A döntéseket precízen meg kell szövegezni (Pl.: Az X objektum Y és Z metódusainak kódját W készíti el Q határidőre).Ha a bejegyzés egyetlen személyhez kötődik, akkor meg kell adni, hogy a tevékenység milyen dologra irányul. A dolog a feladat kapcsán elkészítendő termék, amelynek a (esetleg korábban) beadott anyagban megtalálhatónak kell lenni.A naplóbejegyzés felbontásának egysége szöveges, rajzos anyag esetében az ábra, diagram, vagy kb. fél-egy oldalnyi szöveg. Kódban az egység a metódus. (Pl.: A 3. ábrán látható szekvencia-diagram kidolgozása, vagy az X objektum Y és Z metódusainak kódolása és belövése.
}

\begin{naplo}
    \naplotag{feb. 18. 16h }{ 1 óra }{ Csapat }{ Értekezlet  \newline Döntés: Segédeszközök kiválasztása (git, trello, drive)}
\end{naplo}

\documentclass[../../projlab]{subfiles}
\begin{document}

\makeatletter

\ifSubfilesClassLoaded{
	\coverpage{10. Leadás}
	%\renewcommand{\filePath}[1]{./../../ #1}
	\def\filePath[#1]{./../../#1}
}{}

\makeatother


\chapter{Összefoglalás}
\section{A projektre fordított összes munkaidő}
\comment{Személyenként és összesen}
\begin{ertekelesOra}
    \ertekelestag{Dömötör Péter}{ G2Y5TI }{110 óra, ebből 53 óra önálló és 57 óra csapatmunka}
    \ertekelestag{Sike Ádám}{ E8Z277 }{87 óra, ebből 30 óra önálló és 57 óra csapatmunka }
    \ertekelestag{Gao Tong}{ I2SVOS }{96 óra, ebből 39 óra önálló és 57 óra csapatmunka}
    \ertekelestag{Nagy Beáta}{ GPOGC5 }{98 óra, ebből 41 óra önálló és 57 óra csapatmunka}
    \ertekelestag{Tatai Titusz Miklós}{ IJHLYX }{77 óra, ebből 20 óra önálló és 57 óra csapatmunka}
    \ertekelestag{Csapatgyűlés}{ - }{ 57 óra}
    \ertekelestag{Összesen}{ - }{ $\sum$ }
\end{ertekelesOra}

\section{A feltöltött programok forrássorainak megoszlása}
\begin{ertekelesKod}
    \ertekelestagk{Szkeleton}{3179}
    \ertekelestagk{Prototípus}{5165}
    \ertekelestagk{Grafikus változat}{??}
    \ertekelestagk{Összesen}{ $\sum$  }
\end{ertekelesKod}

\section{Projekt összegzés}
\comment{A projekt tapasztalatait összegző részben a csapatoknak a projektről kialakult véleményét várjuk. A megválaszolandók köre az alábbi}
\begin{itemize}
\item \textit{Mit tanultak a projektből konkrétan és általában?} \newline
    -Fontos, hogy a projekt követelményeit, az elvárásokat pontosan megértsük, mert a félreértésekből rengeteg probléma adódhat később. 
    -A projektmunka megfelelő felépítése és interpretálása is fontos része a folyamatnak. A fejlesztés során rengeteg múlik a kód megfelelő kialakításán. Sokszor számunkra jónak tűnő, működő megoldások során elvi hibákat kaptunk visszajelzésként, így a megfelelő kialakításra a jövőben jobban oda kell figyelnünk. 
    -A csoportmunka hatékony megszervezése nagyon fontos lépés, nagyban meghatározza a projekt sikerességét, és az összes időbefektetést. 
    -A csoportnormák meghatározása, a csoporton belüli megegyezés is fontos tényező, figyelnünk kell rá, hogy olyan kompromisszumokat találjunk (lehetőleg minél hatékonyabban), ami mindenkinek megfelel.  
\item \textit{Mi volt a legnehezebb és a legkönnyebb?} \newline
    -Számunkra a legnehezebb feladatrész a prototípus leadása volt, ezzel töltöttük a legtöbb időt és itt volt szükség sokszor hibajavításra, kódrészletek vagy logika átalakítására. A legkönnyebb beadások pedig az első néhány hét során megalkotott tervek voltak, itt sok ötletünk volt és jó pontszámokat kaptunk. 
    -Általánosságban a határidők betartása, és az időmenedzsment ment nekünk a legnehezebben. A legkönnyebb pedig az ötletelés, a problémákra sokféle megoldás találása volt számunkra.
\item \textit{Összhangban állt-e az idő és a pontszám az elvégzendő feladatokkal?} \newline
    A pontszám összhangban volt a feladatok nehézségével, viszont az idő számunkra nem volt összhangban a feladatok komplexitásával, hosszával. Például volt olyan hét, amikor 10 pontos beadást kellett elkészítenünk, ahol lényegesen kevesebb feladat volt, mint egy 30-40 pontos beadás esetében. Az utóbbi esetben sokkal nehezebb volt jó pontszámot elérnünk, sokkal nagyobb időbefektetés kellett ugyanannyi idő alatt. A heti leadások közti arányon egy kicsit lehetne finomítani, összehangolni, hogy minden héten hasonló nehézségű, hasonló időt igénylő beadást készítsünk. 
\item \textit{Ha nem, akkor hol okozott ez nehézséget?} \newline
    A tárgyon kívül sok másik tárgyra kellett koncentrálnunk hétről hétre, így nem mindig jutott rá elég időnk. A tárgy 3 kreditet ér, és sokszor úgy éreztük, hogy a projekt ennél nagyobb komplexitású, és több időbefektetést igényel. Következő években 3 kreditre egy kicsit kevésbé komplex projektet, vagy több kreditért egy hasonló nagyságrendű projekt kidolgozását ajánljuk. 
\item \textit{Milyen változtatási javaslatuk van?} \newline
    -Sokat segített volna nekünk a leadások során, ha lett volna egy példaprojekt, amihez tartozik minden leadáshoz egy minta dokumentáció. Ezen keresztül pontosabban láttuk volna, hogy a dokumentáció egyes pontjainál helyesen értelmezzük, hogy pontosan miről kellene írnunk. 
    -A kódunk implementációja során sok olyan problémába ütköztünk, ami a tervezés során nem jött elő, nem tudtuk, hogy a gyakorlatban ez nem fog működni. Az UML modellek alkotása során azt tanultuk meg, hogy a tervezésnek minél inkább nyelvfüggetlennek kell lennie, hogy bármilyen környezetben implementálni tudjuk. A gyakorlatban viszont találkoztunk néhány olyan problémával, ami az adott nyelvi környezetben jön elő, és emiatt viszonylag nehéz volt módosítani a modell megfelelő részein. Ezért azt javasoljuk, hogy kicsit nagyobb hangsúly legyen az implementáción a részletes tervezés mellett. 
    -Egy-egy feladatrész (skeleton, prototípus, grafikus verzió) tesztelését manuálisan végeztük, ami során sok hiba nehezen, vagy egyáltalán nem jön elő. Ehelyett a javaslatunk valamilyen tesztosztály, pl. JUnit tesztek elvégzése lenne, ami egy átfogóbb, pontosabb képet adna a kód helyességéről, és jobban kiszűrné a potenciális hibalehetőségeket. 
\item \textit{Milyen feladatot ajánlanának a projektre?} \newline
    -Egy grafikusan kirajzolt sakkjáték (pl. 5D chess w/ multiverse time travel)
    -Egy négyzethálón játszódó kolóniaszimuláció: a pályán emberek mozognak, akikhez feladatokat rendelhetünk (pl. favágás, bányászás), amihez a pályán elhelyezünk alapanyagokat, ezeket kell összeszedniük a tevékenység elvégzéséhez. Bizonyos tevékenységek elvégzése után házakat tudnak építeni maguknak (Rimworld játékhoz hasonló konstrukció)
    -További játékötleteink: Civilization IV, C\&C (régebbiek), Warcraft, Factorio játékok  
\item \textit{Egyéb kritika és javaslat} \newline
    -A tárgyfelelős, Goldschmidt Tanár Úr is tudna esetleg segíteni a projekt sikeres elkészítése érdekében: adhatna ötleteket a megvalósításhoz, írhatna esetleg gyakoribb hibákról és elkerülésükről, arról, hogy mire figyeljünk oda egy-egy beadás során stb. 
    -A projekt során a verziókezelés egy fontos tényező, amit sajnos az eddigi tárgyakban nem tanultunk. Ezért a tárgy első hetén hasznos lenne egy összefoglaló előadás a git használatáról. Elsősorban nem command line megvalósításra gondoltunk, hanem kezdők számára egyszerűbb és átláthatóbb GUI-s megvalósításokra. Ajánljuk a GitKraken, és a SourceTree programokat a verziókezeléshez.
\end{itemize}


\end{document}

\end{document}